\section{Valuations of the finite periods}\label{sec:finite-valuations}
\hypertarget{source-finite-valuations}{}

Section~\ref{sec:transport} identifies the periods for different tuples over the coefficient field. The slope comparison requires lower bounds for their orders at finite places in the prescribed integral coordinates. We estimate the Taylor
coefficients in the local period functions, then transport the
result to the entire column lattice. The prime at which this
divisibility is measured and the auxiliary places used to
realize the row have different roles. The resulting ideal and
index bounds enter the slope identity for the common row in
Section~\ref{sec:word-division}.

\subsection{Two differences of a logarithm}

Two differences of the formal logarithm control the weights of the terms that can first appear in a derivative of the theta kernel.

Fix an odd prime \(p\) and use the valuation \(v(p)=1\) on
\(\C_p\), with \(v(0)=+\infty\). Put
\[
 \rho=(p^2-1)^{-1},\qquad \theta=(p-1)^{-1},\qquad
 \gamma=p\rho,\qquad \theta-\gamma=\rho .
\]
For variables \(X_s\) with input weights \(r_s\), write \(X^n=\prod_s X_s^{n_s}\) for a nonnegative multi-index \(n\). The weighted coefficient valuation of a formal series \(A\) is
\begin{equation}\label{eq:finite-weighted-coefficient-valuation}
 \nu_r(A)=\inf_n\left(v(a_n)+\sum_s r_s n_s\right),
 \qquad A=\sum_n a_n X^n,\qquad \nu_r(0)=+\infty .
\end{equation}
The assertion that \(A\) has weight at least \(h\) means
\(\nu_r(A)\ge h\). Multiplication adds lower bounds on weights. Substitution
preserves a lower bound if the series replacing \(X_s\) has
weight at least \(r_s\) and zero constant term. Both assertions
follow by bounding each summand in each coefficient.

Let \(d\ge2\) be even, index variables by \(\mathbb Z/d\), and
write \(E,O\) for its even and odd classes. For a second copy
of the variables set
\[
 \Delta_E B=B(X_E+Y_E,X_O)-B(X_E,X_O)-B(Y_E,X_O),
\]
and form \(\Delta_O B\) by exchanging the roles of \(E,O\).
On each copy of the variables define
\[
 {\cal L}_s(X)=\sum_{a\ge0}p^{-a}X_{s-2a}^{\,p^{2a}} .
\]
The local arithmetic hypotheses on \(B\) are the following.
Its coefficients lie in a subfield with integral value group.
It is invariant under \(X_s\mapsto\zeta^{(-1)^s p^s}X_s\),
where \(0\le s<d\) and \(\zeta\) has order \(p^d-1\).
Every term has positive degree in each parity block, and the
bidegree \((1,1)\) is absent. For both differences,
\(\exp(\Delta B)({\cal L}(X),{\cal L}(Y))\) has integral
coefficients and constant one. These integrality hypotheses
are on the series in the formal group coordinates.

\begin{lemma}\label{lem:two-differences}
At constant input weight \(\rho\), a coefficient of \(B\) has
weight at least \(\theta\), with the possible exceptions
\[
 X_s^{p^\alpha}X_t^{p^\beta},\qquad
 s\in E,\quad t\in O,\quad s+\alpha=t+\beta\bmod d,\quad
 \alpha+\beta>0 .
\]
Their weight is at least \(\theta-1\).
\end{lemma}
\begin{proof}
The formal inverse \(g\) of \({\cal L}\) exists by induction on
total degree since the linear part is the identity. In
\[
 g_s=X_s-\sum_{a\ge1}p^{-a}g_{s-2a}^{\,p^{2a}},
\]
every coefficient on the right in a nonlinear term uses
lower degrees. Since \((p^{2a}-1)\rho\ge a\), induction gives
weight at least \(\rho\) for each \(g_s\). Apply this substitution
to the integral series in the hypotheses after removing its
constant. We obtain weight at least \(\rho\) for
\(J=\exp(\Delta B)-1\). The torus acts on the second copy with
the same characters as on the first. Hence each nonzero
coefficient of \(J\) is invariant in the sense that its
exponents have total character zero. Modulo \(p+1\), each
character exponent is \(1\), and \(p+1\) divides \(p^d-1\).

For a nonzero coefficient \(j_n=[X^n]J\), the torus invariance
gives \(p+1\mid|n|\). Its valuation is integral because the
additive-coordinate expansion uses rational operations on the
coefficients of \(B\). Therefore
\[
 v_p(j_n)+\rho|n|
     =v_p(j_n)+\theta\,\frac{|n|}{p+1}
     \in \theta\mathbb Z .
\]
The lower bound \(\rho>0\) hence improves to \(\theta\).
Since \(k\theta-v_p(k)\ge\theta\) for every \(k\ge1\), each
term of \(\log(1+J)\) retains the improved bound. This proves
the claim for both \(\Delta_E B\) and \(\Delta_O B\).

Write \(b_n=[X^n]B\). Splitting the even block nontrivially
between its two copies yields in \(\Delta_E B\) the coefficient
\(b_n\prod_{s\in E}\binom{n_s}{m_s}\), where \(m_s\) is the
exponent sent to the second copy. Merging the copies recovers
\(n\) uniquely. If the sum of base-\(p\) digits in one block exceeds one,
send one digit unit \(p^i\) in a coordinate \(s\) to its second
copy. Both parts are nonempty. With digits \(d_j(n_s)\), the
binomial expansion modulo \(p\) gives
\[
   \binom{n_s}{p^i}\equiv d_i(n_s)\not\equiv0\pmod p .
\]

All other binomial factors are one. Thus the weight of \(b_n\)
is at least \(\theta\). Otherwise each block consists of one
\(p\)-power. If, say, \(\alpha>0\), use the split of \(p^\alpha\)
at \(p^{\alpha-1}\), for which
\[
   v_p\binom{p^\alpha}{p^{\alpha-1}}=1 .
\]
The weight is then at least \(\theta-1\). Invariance requires
\(p^{s+\alpha}=p^{t+\beta}\bmod (p^d-1)\).
The order of \(p\) modulo \(p^d-1\) is \(d\), since
\(0<p^i-1<p^d-1\) for \(0<i<d\). This gives the congruence
and also shows that \(\alpha,\beta\) have opposite parity.
\end{proof}

\subsection{The one-letter jet}

Identifying two arguments gives an ordinary derivative in one letter. Its factorial contributes a positive valuation which must be combined with the lower bound for the series coefficients.

For \(d=2\), identify the two letters at each matrix entry
by substituting \((Z_{ib},\eta Z_{ib})\), where \(v(\eta)=0\).
Write \(g_{ib}\) for the inverse of
\(Z\mapsto\sum_{a\ge0}p^{-a}Z^{p^{2a}}\) in that entry.
The one-letter theta series has the form
\[
 \Theta(Z)=U(g(Z))\exp\!\left\{\frac{\pm1}{2}
           \sum_{i,b} B(Z_{ib},\eta Z_{ib})\right\},
 \qquad 1\le i\le N_0,\quad 1\le b\le N_1 ,
\]
with \(U\) integral.

\begin{construction}
\begin{tabular}{@{}l@{\hspace{1.2em}}l@{}}
\toprule
\textit{Contribution to the ordinary derivative} & \textit{Valuation bound} \\
\midrule
Coefficient \([{Z'}^{\kappa}]\Theta(MZ')\) & \(\ge-\gamma|\kappa|\) \\
Multi-index factorial \(\kappa!\) & \(=\theta(|\kappa|-\sum_i s_p(\kappa_i))\) \\
Product, for \(\kappa_i\le a\) & \(\ge\rho|\kappa|-n(1+\log_p a)\) \\
\bottomrule
\end{tabular}
\caption{The valuations entering the one-letter derivative estimate for the integral substitution of Lemma~\ref{lem:one-letter-derivatives}.}\label{fig:analytic-jets}
\end{construction}

\begin{lemma}[Derivatives of the one-letter series]\label{lem:one-letter-derivatives}
Suppose \(B\) satisfies the hypotheses of Lemma~\ref{lem:two-differences}
with \(d=2\). Let \(Z=MZ'\) be an integral linear substitution into
\(n\) variables. For \(a\ge1\) and \(\kappa_i\le a\),
\begin{equation}\label{eq:one-letter-jet}
 v\big(\kappa!\,[{Z'}^\kappa]\Theta(MZ')\big)
 \ \ge\ \rho|\kappa|-n(1+\log_p a).
\end{equation}
\end{lemma}

\begin{proof}
\proofstep{Bound the coefficients before differentiation.}
At input weight \(\gamma\), every term in
the exponent has weight at least \(\theta\). In fact
\(\gamma-\rho=1/(p+1)\), and each exceptional monomial in
Lemma~\ref{lem:two-differences} has degree at least \(p+1\).
The factor \(1/2\) is integral. For the exponential, the term
of order \(k\) has weight at least \(k\theta-v_p(k!)\ge0\).
The integral factor \(U(g(Z))\) also has weight at least zero by
the inverse-series estimate. These conclusions hold after any
integral linear substitution \(Z=MZ'\) into \(n\) variables.
\proofstep{Include the factorials from differentiation.}
The coefficient of \({Z'}^\kappa\) has valuation at least
\(-\gamma|\kappa|\). The factorial identity
\(v_p(m!)=(m-s_p(m))/(p-1)\) therefore gives the lower bound
\[
 (\theta-\gamma)|\kappa|-\theta\sum_i s_p(\kappa_i),
\]
where \(s_p(m)\) is the base-\(p\) digit sum. Each nonzero
\(\kappa_i\) has at most \(1+\log_p a\) digits, each at most
\(p-1\). Since \(\theta-\gamma=\rho\), this is
\eqref{eq:one-letter-jet}.
\end{proof}

\paragraph{Application to framed minors.}

In the column lattice the differential operators are products of framed minors. Their total order specifies which factorial estimate can be used.

For a product of \(C\) minors, let \(h_c\) be their sizes and
\(H=\sum_c h_c\). The corresponding differential operator is
homogeneous of order \(H\). Each entry of the matrix occurs
at most once in a determinant term, hence at most \(C\) times
in a product term.
Thus every derivative monomial \(\partial^\alpha\) in the unframed operator satisfies
\[
  |\alpha|=\sum_{c=1}^C h_c,\qquad 0\le\alpha_{ij}\le C.
\]

At an arbitrary row frame the matrix of
derivations is \(Y\partial_Z\), where the rows of \(Y\) can
be indeterminates. Equivalently apply the unframed minors after
the substitution \(Z=Y^{\mathsf t}Z'\). Use the Gauss valuation
on the coefficient field extended by the entries of \(Y\), under
which the valuation of a polynomial is the minimum of those of
its coefficients. This valuation is multiplicative because
reduction modulo the valuation ideal leaves a polynomial ring
over the residue field, an integral domain. The substitution
matrix is integral, so \eqref{eq:one-letter-jet} gives, for every
coefficient in the frame variables, the same bound
\[
               \rho H-N_0N_1(1+\log_p\max(1,C)).
\]
It also survives any linear reading integral on the frame monomials.
This argument treats the orbit with the two identified letters.
The other even cycles retain a matrix at each letter and use the
individual column heights.

\subsection{Several letters and a linear program}
\label{subsec:cycle-jets}
\hypertarget{source-cyclic-jets}{}

For a full Frobenius orbit the substitutions at successive positions are linked. The exponents contributing to one derivative must therefore satisfy constraints around the entire cycle.

Index the \(d\) matrix variables \(Z_s\) by \(s\in\mathbb Z/d\mathbb Z\). At position \(s\) let the differential operator be a product of minors of the row-framed
matrix \(Y_s\partial_{Z_s}\). A minor of height \(h\) uses the
first \(h\) rows and a specified list of columns. Write \(A_s\)
for the number of minors, \(A'_s\) for the number of height at
least two, and \(H_s\) for the sum of their heights. Choose a set
\(T\subseteq\mathbb Z/d\) of letters at which to treat the first
row separately. Set
\[
 D_s=H_s-{\bf1}_{s\in T}A_s,\qquad
 C_s=\begin{cases} A'_s&s\in T,\\ A_s&s\notin T.\end{cases}
\]
For nonnegative real capacities \(D_s\) allow a variable
\(z_{sj}\ge0\) for each odd \(1\le j<d\).

The variable \(z_{sj}\) uses \(p^jz_{sj}\) units at \(s\) and \(z_{sj}\) at \(s+j\); the sum of all uses
at \(s\) is bounded by \(D_s\). Define
\[
 {\cal F}_p(D)=\gamma\sum_sD_s-\max\sum_{s,j}z_{sj}.
\]
By finite linear programming duality there are numbers
\(0\le b_s\le1\) with
\begin{equation}\label{eq:flow-dual}
 p^jb_s+b_{s+j}\ge1\quad(j>0\text{ odd}),\qquad
 {\cal F}_p(D)=\sum_sD_s(\gamma-b_s).
\end{equation}
Indeed the primal feasible set is nonempty and compact. Its
dual minimizes \(\sum D_s b_s\) over nonnegative \(b_s\)
with the indicated inequalities for \(j<d\). Capping each
\(b_s\) by one preserves feasibility. For any other positive
odd \(j\), reduce \(j\) modulo the even integer \(d\) and use
nonnegativity.

For \(s\in T\) put \(h_s=\min\{j>0:j\ {\rm odd},\ s+j\in T\}\),
with minimum \(\infty\) allowed. Set
\[
 \begin{split}
 T_s={}&\sum_{\substack{0<j<h_s\\ j\ {\rm odd}}}
                 (A_sp^{-j}-C_{s+j})_+
       +\sum_{\substack{j>h_s\\ j\ {\rm odd}}}A_sp^{-j},\\
 W_s={}&\max\{C_{s-j}: j>0\ {\rm odd},\ p^j\le D_{s-j}\},
 \end{split}
\]
where \(t_+=\max(t,0)\) and an empty maximum is zero. For each
destination \(t\) form the list of pairs \((s,j)\), where
\(s\in T\), \(j\) is a signed odd integer, \(s+j=t\), and
\[
       j>0\ \Longrightarrow\ p^j\le A_s,\qquad
       j<0\ \Longrightarrow\ p^{-j}\le D_t .
\]
Denote its length by \(J_t\), and set
\(a_*=\max(1,\max_{s\in T}A_s,\max_t C_t)\).

\begin{proposition}\label{prop:cyclic-jet}
Suppose \(N_0>0\) and \(J_t+{\bf1}_{t\in T}\le N_0\) for all \(t\).
Use \(B\) as in Lemma~\ref{lem:two-differences}, an integral
series \(U\), and the inverse \(g\) of the logarithms \({\cal L}\)
on each matrix entry. Apply the product of all the framed minors
to
\[
          U(g(Z))\exp\!\left(\pm\sum_{i,b}
                        B\big((Z_{sib})_{s\in\mathbb Z/d}\big)\right)
\]
at zero. Each coefficient of the resulting frame polynomial has
valuation at least
\begin{equation}\label{eq:cyclic-gain}
 {\cal F}_p(D)+\sum_{s\in T}(T_s-\gamma W_s)
                    -dN_0N_1(1+\log_p a_*).
\end{equation}
\end{proposition}
\begin{proof}
Use independent indeterminates for the entries of the frames and
the Gauss valuation. At each \(s\in T\) replace the first row
\(Y_{s0}\) by its coordinatewise \(p^m\)-th powers, where
\(m\ge\max_{(s,j)}(-j)\) and \(m\ge0\).
This preserves the valuation of the frame polynomial since
distinct monomials still have distinct exponent vectors. We
construct at \(t\) new rows \(x_{tk}\). The distinguished first
row at \(t\in T\) stays fixed.

To the pair \((s,j)\) ending at
\(t\) assign the covector
\[
                  \lambda_{s,j}=Y_{s0}^{\,p^{m+j}}
\]
using coordinatewise powers of the original first row. Reserve
one distinct nondistinguished row at \(t\) for each of these
covectors. The matrix of their evaluations on the reserved
rows has unit determinant. To see this, reduce the polynomial
modulo the maximal ideal and specialize the reserved rows to
distinct coordinate vectors. These rows are independent of
the covectors, since \(s\ne t\) by parity. In the resulting
determinant each summand is a distinct monomial. At a given
coordinate its variable and exponent recover the pair \((s,j)\)
assigned there. This proves the unit assertion, also when the
list is empty.

Use the inverse evaluation matrix to replace the reserved rows
by dual rows. Subtract the corresponding projections from each
other nondistinguished row. Writing \(r(s,j)\) for the reserved row after this change of basis, the evaluation matrix on all nondistinguished rows becomes
\[
  \lambda_{s,j}(x_{tk})=\begin{cases}1&k=r(s,j),\\0&k\ne r(s,j).\end{cases}
\]
Both basis changes have integral coefficients.
The old minor operators expand in these new rows with integral
coefficients. A determinant with two occurrences of the same
row vanishes.

In any surviving monomial of differentiation with
multi-index \(\kappa\), the distinguished row has degree exactly
\(A_s\), the other rows at \(s\) total \(D_s\), and each of them
has degree at most \(C_s\). At a selected letter, only a minor
of height at least two can contribute to a nondistinguished row;
the row change mixed only those rows with each other. In
particular every coordinate of \(\kappa\) is at most \(a_*\).
The row-by-row bounds on the differentiation orders read
\[
  \sum_b\kappa_{s0b}=A_s,\qquad \sum_{k\ne0,b}\kappa_{skb}=D_s,\qquad \sum_b\kappa_{skb}\le C_s\quad(k\ne0).
\]

In evaluating this monomial make the dual change of variables
\(Z_{sab}=\sum_k x_{sk,a}Z'_{skb}\). It suffices to bound
the series modulo the monomial ideal generated by
\(({Z'}_i^{\kappa_i+1})_i\), since all coefficients entering
\([{Z'}^\kappa]\) survive there. Initially give every variable
weight \(\rho\). By Lemma~\ref{lem:two-differences} only the
exceptional terms need an extra unit of weight for the
exponential. In the expansion of such a term, a mixed choice of
coordinates in either \(p\)-power of a linear form has
multinomial coefficient divisible by \(p\). For the remaining
choices write \(s+\alpha=t+\beta\), with \(s,t\) distinct and of
opposite parity. For coordinates \(i=(s,k,b)\), \(i'=(t,k',b')\)
the bad product is possible only for \(b=b'\). Its coefficient
from that term, divided by the original coefficient of \(B\),
is
\[
                \sum_a x_{sk,a}^{\,p^\alpha}
                           x_{tk',a}^{\,p^\beta}.
\]

If \(i\) is distinguished and \(i'\) is not, let
\(j=\alpha-\beta\). If the product survives truncation then
\((s,j)\) belongs to the list at \(t\). If \(k'\) is not its
reserved row, the last display is zero modulo \(p\). For \(j>0\)
it is congruent to
\((\lambda_{s,j}(x_{tk'}))^{p^\beta}\), and for \(j<0\) it is
congruent to the same evaluation raised to \(p^\beta\), using
\(x_{s0}=Y_{s0}^{p^m}\). Both congruences hold modulo the ideal
\((p)\) of the integral coefficient ring by the binomial
theorem.
In particular, for an unreserved partner the evaluation contributes
\[
  \lambda_{s,j}(x_{tk'})=0\quad\Longrightarrow\quad \sum_a x_{s0,a}^{p^\alpha}x_{tk',a}^{p^\beta}\in(p).
\]
 Thus this case also recovers the missing unit.

We specify extra input weights \(r_i\) for the remaining
products. Give a distinguished row at \(s\) weight increment
\(\epsilon_s\ge0\), requiring \(\epsilon_s\ge p^{-h_s}\) for
finite \(h_s\). Give the nondistinguished rows at \(t\) increment
\(b_t\), raised on the row reserved for \((s,j)\) to
\[
 \begin{cases}
 \max(b_t,1-p^j\epsilon_s),&j>0,\\
 \max(b_t,p^j),&j<0 .
 \end{cases}
\]
For two nondistinguished rows, \eqref{eq:flow-dual} gives
\(p^\alpha r_i+p^\beta r_{i'}\ge1\) after orienting the difference
of exponents so that it is positive.

For two distinguished rows
in that orientation, the first has \(h_s\le\alpha-\beta\), so
the same bound follows from its own increment. For one row of
each kind, the reserved-row formula gives the bound with either
sign of \(j\). Every surviving exponent term now has weight at
least \(\theta\) at input weights \(\rho+r_i\). The exponential
and the integral factor therefore have weight at least zero.
Using the factorial formula as in \eqref{eq:one-letter-jet}, we
obtain the lower bound
\[
       \sum_i\kappa_i(\gamma-r_i)-dN_0N_1(1+\log_p a_*).
\]

For a fixed \(s\in T\) the extra cost of the negative labels
is at most \(W_s\sum_{j>0,\ j\ {\rm odd}}p^{-j}=\gamma W_s\).
For the positive labels and the distinguished row, one can choose
\(\epsilon=\epsilon_s\) so that
\[
 A_s\epsilon+
 \sum_{\substack{j>0\ {\rm odd}\\(s,j)\ {\rm in\ a\ list}}}
       C_{s+j}(1-p^j\epsilon)_+
                \ \le\ A_s\gamma-T_s .
\]
Indeed the right side is the sum of
\(\min(A_sp^{-j},C_{s+j})\) at odd \(j<h_s\), together with
\(A_sp^{-h_s}\) when \(h_s\) is finite. Take the first odd
\(j<h_s\) with \(A_sp^{-j}\le C_{s+j}\), or \(j=h_s\) if none
exists.

At a finite choice set \(\epsilon=p^{-j}\).
For a finite choice the threshold has the two consequences
\[
  A_sp^{-i}>C_{s+i}\quad(0<i<j,\ i\ {\rm odd}),\qquad p^i\epsilon\ge1\quad(i\ge j).
\]
 All earlier
positive-label costs are bounded by their \(C_{s+i}\); later
ones and the one at \(j\) cost zero. The term at \(j\) pays for
\(A_s\epsilon\). If the choice is infinite use \(\epsilon=0\)
and all finite costs are bounded by the respective minima.
This also verifies the lower bound on \(\epsilon_s\). The
nondistinguished baseline is
\(\sum D_t(\gamma-b_t)={\cal F}_p(D)\), and the distinguished
baseline at \(s\), before paying costs, is \(\gamma A_s\).
Adding gives \eqref{eq:cyclic-gain}.
\end{proof}

For comparison with the common chain, write \(V_s=(k-1)w_s\)
with \(k\ge3,\ w_s>0\), and suppose
\[
      D_s=V_s+E_s-2{\bf1}_{s\in T}w_s,\qquad
                         0\le E_s\le4w_s .
\]
Then
\begin{equation}\label{eq:flow-comparison}
            {\cal F}_p(D)\ge {\cal F}_p(V)
                                   -\frac{11}{p}\sum_s w_s .
\end{equation}
To prove this take an optimum \(b\) in \eqref{eq:flow-dual}
for \(V\), and let \(H=\{s:b_s\ge2/p\}\). Complementary
slackness with an optimum flow shows that every vertex of \(H\)
has saturated capacity, and every edge starting there has zero
flow, since its dual constraint is strict. Hence
\(\sum_{s\in H}V_s\) is the total flow into \(H\). Summing the
capacity inequalities shows \(p\sum z_{sj}\le\sum_sV_s\), so

\[
  \sum_H w_s\le p^{-1}\sum_s w_s.
\]
 In the dual formula,
increasing \(V\) by \(E\) decreases its objective at this \(b\)
by at most \(\sum E_s(b_s-\gamma)_+\). On \(H\) use \(b_s\le1\);
outside \(H\) use \(\gamma\ge1/p\). The resulting loss is at
most \((8/p)\sum w_s\). Finally adding \(2{\bf1}_T w\) to
\(D\) increases its optimum dual objective by at most
\(2\gamma\sum_T w_s\), since every feasible \(b\) is
nonnegative. As \(2\gamma\le3/p\), this proves
\eqref{eq:flow-comparison}.
Together the two changes of capacities cost at most
\[
  \frac8p\sum_s w_s+2\gamma\sum_s w_s\le\frac{11}{p}\sum_s w_s.
\]

Taking the smaller numerator \(A_s-A'_s\) and subtracting
\(C_{s+j}\) at every odd \(j\ne h_s\) can only reduce \(T_s\).
Replacing \(D_{s-j}\) by \(H_{s-j}\) in the definition of \(W_s\)
can only increase the loss. If \(H_s\le a\), the number of
positive and negative labels at one destination is at most
\(2\log_p a+2\): for each signed \(j\) there is just one source.
Thus \(2\log_p a+2\le N_0-1\) suffices for the row count; also
\(a_*\le\max(1,a)\). These observations give the constants
\(11\) for the capacity error and \(1\) for the digit loss in
the split bound. A common unit rescaling of the local variables
multiplies a homogeneous jet by a unit and leaves its valuation
unchanged.  \label{cor:cyclic-jets-1}

\subsection{Evaluating a derivative on an integral lattice}
\hypertarget{source-integral-jets}{}

The coefficient estimates so far use products of framed minors. To pass to integral vectors of their rational span one needs saturation over the local valuation ring.

\begin{lemma}[Saturation of flag-minor polynomials]\label{lem:saturated-minor-polynomials}
Over a valuation ring \(\mathcal O\), fix a list of column
heights and take products of minors in the first \(h\) rows
at each height \(h\). Write \(S\) for their \(\mathcal O\)-span and \(A\) for the polynomial ring. For every \(f\in A\) and \(a\in\mathcal O\),
\[
 a\ne0,\quad af\in S\quad\Longrightarrow\quad f\in S.
\]

\end{lemma}

\begin{proof}
Sort the indices within each tableau column,
carrying the determinant sign; a repeated index kills the term.
One may straighten until the indices across each tableau row
are weakly increasing. Indeed suppose adjacent columns
\(a_1<\cdots<a_h\), \(b_1<\cdots<b_k\), \(k\le h\),
violate \(a_i\le b_i\). Antisymmetrize the two-minor product
in the \(h+1\) vectors indexed by
\(a_i,\ldots,a_h,b_1,\ldots,b_i\).

The antisymmetrized product vanishes because both minors use only the first \(h\) coordinates. Grouping shuffles
modulo permutations internal to the two lists gives an identity
with coefficients \(\pm1\) before collecting like terms. This
identity holds over the integers, as is seen by first expanding
over \(\mathbb Q\) and cancelling the common factorials.
With \(P\) the identity-shuffle product, the alternating relation has the form
\[
  P=-\sum_{\sigma\ne1}\operatorname{sgn}(\sigma)P_\sigma,
\]
where the sum uses the shuffle representatives and each \(P_\sigma\) retains the order inside the two lists. Every other shuffle
moves a value at most \(b_i<a_i\) into the left column in place
of a larger value. After sorting, a surviving term has a
strictly smaller left column in lexicographic order. Earlier
columns remain fixed. The whole tableau decreases in a finite
order, which proves spanning.

Order the matrix variables row by row, preferring smaller
column indices, and take lexicographic monomial order. The
leading monomial of a sorted minor is its diagonal product
with coefficient one. In a straightened tableau the row
multisets read off from the diagonal monomial recover each
weakly increasing row.
For a tableau with column entries \(a_{1c}<\cdots<a_{h_c c}\), the recovery rule is
\[
  \operatorname{in}(P)=\prod_c\prod_{i=1}^{h_c}X_{i,a_{ic}},\qquad \deg_{X_{ij}}\operatorname{in}(P)=\#\{c:a_{ic}=j\}.
\]
 Distinct such tableaux consequently have
distinct leading monomials. In an expansion over the fraction
field, the coefficient of the largest leading monomial is the
coefficient of its tableau. Subtract it and continue downwards;
if the polynomial is integral all coefficients so obtained
are integral. This proof works for products at several
letters because their matrices have disjoint variables.

\end{proof}

\begin{lemma}[Transfer of a coefficient bound to integral jet functionals]
\label{lem:jet-functional-coefficient-bound}
If \(\lambda\) is integral on the prescribed products of flag minors and \(P=\sum_\alpha c_\alpha X^\alpha\) is a local jet polynomial, then
\[
 v(\lambda(P))\ge\min_\alpha v(c_\alpha).
\]
 Joint integral remainders in independent
orbit variables obey the sum of the individual orbit bounds.
\end{lemma}

\begin{proof}
Write the jet polynomial in the frame \(Y_s\).
For a minor with row set \(I\) and column set \(J\),
Cauchy--Binet applies because the constant-coefficient
derivations commute:
\[
    \det((Y_s\partial)_{I,J})
      =\sum_{|K|=|I|}\det((Y_s)_{I,K})\det(\partial_{K,J}).
\]
The row set \(I\) is initial. After those derivatives act and
are evaluated at zero, only scalar coefficients of products
of initial-row minors remain. These are the polynomials of
Lemma~\ref{lem:saturated-minor-polynomials}. If the resulting
polynomial is nonzero with coefficient of least valuation \(a\),
divide by that coefficient to obtain an integral polynomial
in their fraction-field span. Saturation gives an integral
linear combination. Applying the integral functional to this
combination gives a value of valuation at least \(v(a)\)
after scaling back. The zero polynomial contributes zero.

For independent
sets of orbit variables an arbitrary integral factor \(U\) may
depend on all sets. Expand it in monomials before applying the
substitutions \(g\). Each monomial factors over the orbits, with
integral scalar, and only finitely many affect a specified jet
since \(g(0)=0\). Hence the bounds add even for this joint \(U\).
\end{proof}

There is a different calculation for a ramified quadratic factor.
At a place over \(p\), suppose the local series has the form
\[
 \exp\!\left(-S\sum_{t,a,b}u_t X_{tab}Y_{tab}\right) H(X,Y),
 \qquad v(S)\ge-1/2,\quad v(u_t)\ge0 ,
\]
and that every derivative of \(H\) at zero is integral. This last
property is preserved by integral linear substitutions, by the
chain rule with constant coefficients. Let \(H^\pm_t,A^\pm_t\)
be the total orders and numbers of minors on the two sides,
with \(X\) the minus side, and let \((A^-_t)'\) count minus
minors of height at least two.

For a fixed differential monomial
with respective entry exponents \(m^\pm_{tab}\), in the
derivative of the exponential times \(H\) each factor \(S\)
uses one differentiation from each side at a common entry.
At one entry write \(m^\pm=m^\pm_{tab}\). Choose \(h\) derivatives
from each side and match them by a bijection. The corresponding
coefficient and its valuation satisfy
\[
 C_h=\binom{m^-}{h}\binom{m^+}{h}h!(-Su_t)^h,
 \qquad v(C_h)\ge -h/2,\qquad 0\le h\le\min(m^-,m^+).
\]
The remaining derivatives act on \(H\) and have integral value. Summing
the losses over entries gives the lower bound
\(-\frac12\sum_{t,a,b}\min(m^-_{tab},m^+_{tab})\).

Evaluate the frames at independent indeterminates with their
Gauss valuation. The determinants of the square frames are
units there. Choose as the new minus derivations its given rows
and as the new plus derivations the inverse transpose rows.
For a change \(U\) of minus frame, the inverse transpose on the plus frame preserves the contraction because
\[
  (U x)^{\mathsf t}(U^{-\mathsf t}y)=x^{\mathsf t}y.
\]
 The plus frame is now an integral combination of the
new plus rows, and the minus minor operators use the identity
frame. Without restriction the sum of the minima at \(t\)
is at most \(H_t^-\).

At a selected \(t\in T\), the minus
first-row degree is exactly \(A^-_t\), while the plus first-row
degree is at most \(A^+_t\) after expanding its minors. Use
the plus count on that row and the minus counts on the other
rows. Writing \(a=1\) for the first row, this gives
\[
  \sum_{a,b}\min(m^-_{tab},m^+_{tab})\le \sum_b m^+_{t1b}+\sum_{a\ne1,b}m^-_{tab}\le A_t^+ +H_t^- - A_t^-.
\]
In particular the uniform lower bound used in the chain is
\[
 -\frac12\sum_t\left[H_t^-
       -{\bf1}_{t\in T}\{A^-_t-(A^-_t)'\}
       +{\bf1}_{t\in T} A^+_t\right].
\]
Passing back to the generic frame gives this bound coefficientwise.
An integral functional on its monomials preserves it. At a
place where all derivatives are integral and the differential
operator has integral coefficients, the bound is simply zero.

The calculations apply when the corresponding local form and
integral remainder have been obtained from the curve.
\label{cor:integral-jets-1}

\subsection{The exponent from an elliptic curve}
\label{subsec:curve-exponents}
\hypertarget{source-curve-exponents}{}

The logarithmic differences and torus action used in the jet estimates arise from a single elliptic formal group. We construct them in coordinates in which the Frobenius congruences can be compared. Work first over a finite
unramified extension of \(\mathbb Q_p\), or the completion of its maximal
unramified extension, with Frobenius lift \(\sigma\), \(p\ge5\).
Write \(f\sim h\) for zero-constant series if \(\exp(f-h)\)
is an integral unit with constant one. The relation is additive,
is preserved by integral substitutions fixing zero, and permits
multiplication by a scalar in \(\mathbb Z_p\). For the latter
property expand a principal unit to that power by the integral
binomial coefficients. These statements continue to hold after
valued extension.

Let 
\[
  L(X)=\sum_{n\ge0}p^{-n}X^{p^{2n}},
\]
 and let \(F\) be
the formal law with logarithm \(L\). We need precise control
of which additive-coordinate substitutions are integral in
formal group coordinates. On series in any finite set of
variables let \(\Phi\) apply \(\sigma^2\) to the coefficients
and raise the variables to \(p^2\). For \(A(0)=0\),
\begin{equation}\label{eq:integral-log}
 L^{-1}(A)\text{ is integral}\quad\Longleftrightarrow\quad
                       A-p^{-1}\Phi A\text{ is integral}.
\end{equation}
\begin{proposition}[The integral logarithm criterion]\label{prop:quadratic-integral-log}
On the unramified ring and for \(A(0)=0\), the condition
\eqref{eq:integral-log} characterizes integral substitutions into
\(L^{-1}\). In particular the following formal sum and scalar substitutions have coefficients in the unramified ring, with \(a\) integral in the scalar case:
\[
 F(X,Y)=L^{-1}(L(X)+L(Y)),\qquad [a]X=L^{-1}(aL(X))\quad(\sigma^2a=a).
\]

\end{proposition}

\begin{proof}
For \(A=L(h)\) the latter difference is \(h\) plus
\[
 \sum_{a\ge0}p^{-a-1}\big\{(h^{p^2})^{p^{2a}}-
                                      (\Phi h)^{p^{2a}}\big\}.
\]
If \(h\) is integral, then \(h^{p^2}\equiv\Phi h\bmod p\).
Raising a congruence modulo \(p^b,\ b\ge1\), to the \(p\)-th
power improves it to one modulo \(p^{b+1}\). This proves
integrality termwise. Conversely determine \(h\) by degree:
at each degree the terms in braces depend only on earlier
coefficients, so the same argument applied to that integral
truncation proves the next coefficient integral. In particular
\(F\) and \([a]X=L^{-1}(aL(X))\) for \(a\) integral,
\(\sigma^2a=a\), are integral.

\end{proof}

Use a short Weierstrass curve \(y^2=4x^3-g_2x-g_3\) with
integral invariants and good reduction. Let \(z\) be its
additive uniformizing parameter, \(x=\wp(z),y=\wp'(z)\),
and \(t=-2x/y\) the formal parameter. Then \(z=l(t)=t+\cdots\)
is its formal logarithm; \(x(t)=t^{-2}a(t)\) with \(a\)
integral, even, \(a(0)=1\). Indeed, substituting \(y=-2x/t\) and \(x=t^{-2}a\) in the curve
equation gives
\[
 P(t,a)=a^3-a^2-(g_2/4)at^4-(g_3/4)t^6=0,\qquad
 P(0,1)=0,\quad \partial_aP(0,1)=3-2=1 .
\]
The formal implicit-function theorem therefore gives a unique integral
solution with \(a(0)=1\).

Since \(P(-t,a)=P(t,a)\), uniqueness also
gives \(a(-t)=a(t)\). Write \(\varsigma\) for
the odd Weierstrass sigma function normalized by
\(\varsigma'(0)=1\). Its formal series is specified by
\(-(\log\varsigma)''=\wp\) and oddness with the usual
zero quadratic term in \(\log(\varsigma(z)/z)\).
For a scalar \(s_2\) put
\[
       q_E(t)=-\log\frac{\varsigma(l(t))}{t}
                                     +\tfrac12s_2 l(t)^2.
\]
Let \(t_\pm\) represent addition and subtraction on the curve.
\begin{proposition}[Additivity and conjugate scalars in the curve pairing]\label{prop:curve-additive-pairing}
Suppose the good-model change \(l^{-1}(cL(X))\) is integral
with \(c\) a unit. The parallelogram relation for \(q\) holds
modulo \(\sim\), and its polarization \(b\) is additive in both
variables modulo \(\sim\). If \eqref{eq:cm-norm-relation} holds
for all quadratic units, then it holds for all quadratic integers
and for each such \(a\) the polarized scalar action satisfies
\[
 b([a]x,y)\sim b(x,[\sigma(a)]y).
\]

\end{proposition}

\begin{proof}
\proofstep{The integral parallelogram relation.}

The formal addition formula for sigma gives
\[
 \frac{\varsigma(l(t)+l(u))\varsigma(l(t)-l(u))}
      {\varsigma(l(t))^2\varsigma(l(u))^2}=\wp(l(u))-\wp(l(t)).
\]
For instance, as elliptic functions of \(l(t)\) the two sides
have the same poles and principal parts by quasi-periodicity
and the normalization at zero. Their difference is constant,
and evaluation at \(l(t)=l(u)\) fixes that constant. Expansion
at the identity yields the formal identity over the field of
the invariants by the coefficient recursions for \(\wp\) and
\(\varsigma(z)/z\).

It follows that the exponential of minus the parallelogram
difference \(q_E(t_+)+q_E(t_-)-2q_E(t)-2q_E(u)\) is
\[
 \frac{t^2a(u)-u^2a(t)}{t^2-u^2}
                   \frac{t+u}{t_+}\frac{t-u}{t_-}.
\]
Each factor is an integral unit with constant one. Evenness
of \(a\) gives this for the first; integrality of the curve's
formal law and its odd inverse give divisibility by \(t\pm u\)
and the linear term gives the other two. For a unit \(c\)
suppose that \(t(X)=l^{-1}(cL(X))\) is integral. This is the
coordinate comparison required of the local CM model. Its
linear term is \(cX\). The exponent in these coordinates is
\[
           q(X)=-\log\frac{\varsigma(cL(X))}{cX}
                                    +\tfrac12s_2(cL(X))^2 .
\]
It differs from \(q_E(t(X))\) by a logarithm of an integral unit
since \(t(X)/(cX)\) is such a unit. The same parallelogram
relation therefore holds for \(q\) modulo \(\sim\).

\proofstep{Additivity via the integral formal half.}

Let \(b(X,Y)=q(X+_FY)-q(X)-q(Y)\). It vanishes on the two
axes and is symmetric. Modulo \(\sim\) it is additive in each
variable on all integral formal points. To see this, subtract
the parallelogram identities at \((x+z,y)\) and \((x,y)\),
writing the group law additively. For \(f(x)=b(x,z)\) this gives
\(f(x+y)+f(x-y)\sim2f(x)\). Taking \(y=x\) gives

\[
  f(2x)\sim2f(x)-f(0)=2f(x),
\]
 since \(f(0)=0\). With \(x=(u+v)/2\) and \(y=(u-v)/2\) in the formal group, the two
relations give
\[
 f(u)+f(v)\sim 2f\bigl((u+v)/2\bigr)\sim f(u+v).
\]
The formal half is integral here. Symmetry of \(b\) gives additivity
in the second variable.

\proofstep{Conjugate scalars and extension from units.}

There is also a scalar comparison from CM. Its needed input
on this model is that for every unit \(a\) fixed by \(\sigma^2\),
\begin{equation}\label{eq:cm-norm-relation}
                   q([a]X)\sim a\sigma(a)\,q(X).
\end{equation}
The scalar \(n=a\sigma(a)\) is in \(\mathbb Z_p^\times\).
Polarization gives \(b([a]x,[a]y)\sim n b(x,y)\).
For \(n\in\mathbb Z_p\), additivity implies
\(b(x,[n]y)\sim n b(x,y)\): approximate \(n\) by nonnegative
integers. Each coefficient of each side and its exponential
difference is continuous in \(n\), since at bounded degree
only finitely many operations occur. The valuation ring is
closed. Replacing \(y\) by \([a^{-1}]y\) therefore gives
\[
                        b([a]x,y)\sim b(x,[\sigma(a)]y).
\]
This holds for nonunits \(a\) fixed by \(\sigma^2\) as well,
by subtracting the identities for \(1+a\) and \(1\).

In particular \(q([a]x)\sim a\sigma(a)q(x)\) for nonunits too:
use \([a]x=[1+a]x-x\), the evenness of \(q\),
and \(b(x,x)\sim2q(x)\) from its parallelogram identity.
Moving \(1+a\) across this pairing, or polarizing
\(q([1+a]x)\), gives the claimed identity. Explicitly
\(b([a]x,x)+b(x,[a]x)\sim
b(x,[a+\sigma(a)]x)\sim2(a+\sigma(a))q(x)\);
this determines the cross term by division by two.

\end{proof}

Write \(\widetilde b(G,H)=b(L^{-1}G,L^{-1}H)\).
The preceding rules apply whenever \(G,H\) satisfy
\eqref{eq:integral-log}. For an even cycle choose integral
bases \(\omega_r,\omega_r^*\) of the unramified degree \(d\)
field over its unramified quadratic subfield, dual for the
relative trace. Such bases exist since the residue trace is
perfect and its lifted determinant is a unit. Set
\[
 \begin{aligned}
 E(c,Z)&=\sum_{s\ {\rm even}}\sigma^s(c)Z_s,&
 O(c,Z)&=\sum_{s\ {\rm odd}}\sigma^s(c)Z_s,\\
 B^\circ(Z)&=\sum_r\widetilde b(E(\omega_r^*,Z),O(\omega_r,Z)).
 \end{aligned}
\]
Define \(B\) as the average of \(B^\circ\) under
\(Z_s\mapsto u^{(-1)^s p^s}Z_s\) over \(u^{p^d-1}=1\). It
has coefficients in an unramified field and is torus invariant.
The two positive parity degrees follow from the axes of \(b\).
An even-odd bilinear term would have character
\(p^s-p^t\bmod(p^d-1)\) for distinct \(s,t\), so the averaging
eliminates it.

\begin{proposition}[The cycle exponent in a trace-dual frame]\label{prop:curve-cycle-exponent}
On an even cycle, with the integral curve and scalar comparisons
just proved, the averaged exponent \(B\) has the two integrality
properties of Lemma~\ref{lem:two-differences}. The trace-dual
substitution from an integral orthogonal theta chart gives
the local series of Proposition~\ref{prop:cyclic-jet}, whose exponential factor and integral remainder combine to give
\[
 \exp\!\left(-\sum_{i,b}B((Z_{sib})_s)\right)U(g(Z)).
\]
 The one-letter exponent is the equal-variable
specialization with the factor \(1/2\).
\end{proposition}

\begin{proof}
\proofstep{Check the two parity differences.}

For \(Z_s={\cal L}_s(X)\) one has
\(Z_s-p^{-1}\Phi Z_{s-2}=X_s\). More generally suppose all
these differences are integral and all \(Z_s(0)=0\). Both
\(E(c,Z)\) and \(O(c,Z)\) satisfy the right side of
\eqref{eq:integral-log} for integral \(c\) fixed by \(\sigma^d\),
by shifting indices in the sum. Torus translation replaces
\(\omega_r^*\) by \(u\omega_r^*\) and \(\omega_r\) by
\(u^{-1}\omega_r\), since 
\[
  \sigma^s(u)=u^{p^s}.
\]
The parity differences of each translated \(B^\circ\), evaluated
on the two copies of \({\cal L}\), are therefore equivalent to
zero by additivity. This remains true on averaging since
\((p^d-1)^{-1}\in\mathbb Z_p\). This verifies precisely the two
integrality hypotheses of Lemma~\ref{lem:two-differences}.

\proofstep{Compare the torus average to the original contraction.}

To identify the resulting exponent with the one on the original
theta chart, one needs \(B(Z)\sim B^\circ(Z)\) on these points.
For a torus element write
\(c_{rr'}=\operatorname{Tr}(u\omega_r^*\omega_{r'})\), with
the trace over the quadratic subfield. Its entries are integral,
fixed by \(\sigma^2\). Then
\[
 u\omega_r^*=\sum_{r'}c_{rr'}\omega_{r'}^*,\qquad
 \omega_{r'}=\sum_r c_{rr'}\,u^{-1}\omega_r .
\]
In \(E\) the coefficients of the first expansion are \(c_{rr'}\).
Move each across \(\widetilde b\) by the scalar comparison;
the resulting \(\sigma(c_{rr'})\) are exactly the coefficients
in \(O\) for the second expansion. This proves invariance
modulo \(\sim\) under each torus element, whence the asserted
comparison with the mean.

\proofstep{Apply the hyperbolic Gram identities.}

For clarity, suppose the theta series on the orthogonal chart
has \(U_0=\exp(\sum_i q(X_i))\Theta(X)\) integral. To pass to
a hyperbolic frame use the formal sum of multiples of independent
coordinates with coefficients \(f_{ei},f'_{ei}\) integral in
the quadratic subfield. The orthogonality conditions on these two families are
\[
  \sum_i\sigma(f_{ei})f_{e'i}=0,\qquad\sum_i\sigma(f'_{ei})f'_{e'i}=0,\qquad\sum_i\sigma(f_{ei})f'_{e'i}=\delta_{ee'}.
\]
Expanding \(q\) on this formal sum gives diagonal terms
\(f_{ei}\sigma(f_{ei})q(X_e)\) and terms paired by \(b\),
with the conjugated first coefficient moved to the second slot.
Summing over \(i\) leaves just \(\sum_e b(X_e,Y_e)\) modulo
\(\sim\). Thus after this substitution the theta series is an
integral series times the exponential of minus that sum.
The index \(e\) includes a trace-basis index and a matrix
entry.

Now set
\[
  X_r=L^{-1}(E(\omega_r^*,Z)),\qquad
  Y_r=L^{-1}(O(\omega_r,Z))
\]
at each entry. On substituting \(Z={\cal L}(W)\), both are
integral formal points by \eqref{eq:integral-log}. The
exponential of the difference between the summed exponents
and the sum of the corresponding \(B\)'s is an integral unit
in \(W\). Absorb it into the integral factor and substitute
\(W=g(Z)\). This gives the factor \(U(g(Z))\) and the exponent
\(-\sum_{i,b} B((Z_{sib})_s)\) of
Proposition~\ref{prop:cyclic-jet}. Auxiliary orbit variables
can be left as indeterminates throughout. For one letter use
\(d=2,\ \omega=\omega^*=1\). In the equal-variable
specialization 
\[
  b(X,X)=q([2]X)-2q(X)\sim 2q(X);
\]
 division by two yields
the one-letter exponent. A unit change of the logarithmic
variable common to the two slots scales the homogeneous jet by
a unit.

This calculation gives the local series from an integral
orthogonal theta chart, an integral logarithm comparison and
\eqref{eq:cm-norm-relation}. Establishing those three inputs
on the family requires its models and the normalization of its
translated theta sections.
\label{cor:curve-exponents-1}
\end{proof}

\subsection{Good inert places and translated theta sections}
\label{subsec:cm-models}
\hypertarget{source-cm-models}{}

The formal group computation applies to the period kernel only if the translated theta sections extend on the corresponding integral CM model.

Let \(K=\mathbb Q(\sqrt{-D})\), where \(D>0\) is squarefree
and \(D\equiv3\bmod4\). The period lattice \(\Gamma\subset\C\)
is stable under \((1+\sqrt{-D})/2\). Thus its multiplier order
is the full ring of integers of \(K\). Fix \(p\ge5\) inert in
\(K\) and an abstract field isomorphism \(\C\simeq\C_p\).
All analytic identities below are transported on Taylor
coefficients. Put \(R=W(\overline{\mathbb F}_p)\), the integers
of the completed maximal unramified field, and let \(\sigma\)
be its Witt Frobenius.

\begin{proposition}[Good model and integral CM coordinates]\label{prop:good-inert-cm-model}
At the specified inert \(p\ge5\), a homothety of the complex
period lattice gives a good short model with coefficients in \(R\)
under the chosen transport. The
elliptic formal group has \(p\)-height two and admits an integral comparison with unit tangent
\[
 t(X)=l^{-1}(cL(X))\in R[\![X]\!],\qquad t'(0)=c\in R^\times.
\]

Nonzero prime-to-\(p\) torsion sections have integral affine
coordinates on the good model.
\end{proposition}

\begin{proof}
\proofstep{Obtain good reduction with full level three.}

Equip the elliptic curve with a
full level-three basis. The main theorem of complex multiplication
gives the moduli point with this basis over the ray class field
of modulus \(3\) over \(K\): under reciprocity the local units
preserve the ideal lattice and the units congruent to \(1\) at
\(3\) preserve the basis modulo \(3\). The full basis removes
automorphisms, so the curve with level structure descends there.
This field is unramified at \(p\). The singular modulus is
integral, so the curve has potentially good reduction. Its
inertia image on the \(3\)-adic Tate module is finite and
congruent to the identity modulo \(3\); it is trivial, since
the matrix logarithm on \(1+3M_2(\mathbb Z_3)\) has the
matrix exponential as inverse, and a torsion matrix there
has logarithm zero.
Indeed a torsion element \(A\in1+3M_2(\mathbb Z_3)\) of order \(n\) satisfies
\[
  0=\log(A^n)=n\log A\quad\Longrightarrow\quad \log A=0\quad\Longrightarrow\quad A=I .
\]

The good reduction criterion applies
since \(3\ne p\). A good Weierstrass model can be shortened
over \(R\) because \(2,3\) are units. Matching an invariant
differential scales \(g_2,g_3\) by the fourth and sixth
powers of a single factor.
Thus a change of scale by \(u\) preserves the quotient defining \(j\), since
\[
  (u^4g_2)^3=u^{12}g_2^3,\qquad (u^6g_3)^2=u^{12}g_3^2.
\]
 CM endomorphisms extend over the
valuation ring by the extension property of abelian schemes;
they are defined over the moduli field containing \(K\),
as can also be seen from their faithful action on the tangent
line. \cite{cm-reciprocity,serre-tate}

\proofstep{Determine the height and lift the formal isomorphism.}

The formal completion is a one-dimensional formal
\({\cal O}_{K_p}\)-module. The action of \({\cal O}_K\) extends
by \(p\)-adic continuity: in logarithmic coordinates the
coefficients of \([a](t)=l^{-1}(a l(t))\) at each degree are
polynomials in \(a\), and \(R\) is closed. The special fibre
of the elliptic formal group has \(p\)-height \(h=1\) or \(2\).
The first nonzero term of \([p]\) in the residue field has
degree \(p^h\). Let its leading coefficient be \(A\ne0\).

In the residue
field the leading terms of \([p]\circ[a]\) and
\([a]\circ[p]\) agree, so
\[
                    A\bar a^{p^h}= \bar a A
                            \quad(\bar a\in\mathbb F_{p^2}).
\]
Thus \(2\mid h\), and \(h\in\{1,2\}\) gives \(h=2\). The height as
a formal \({\cal O}_{K_p}\)-module is one. Over an algebraically
closed residue field such modules of the same height and
structural embedding are isomorphic, and the deformation ring
of a height-one formal module is just \(R\). The model with
logarithm \(L\) in \eqref{eq:integral-log} has this height:
its \([p]\)-series has linear coefficient \(p\), no terms of
degrees \(2,\ldots,p^2-1\), and coefficient \(1-p^{p^2-1}\)
at degree \(p^2\), by comparison in \(L([p]X)=pL(X)\).

In particular the first surviving term on reduction is
\[
  [p](X)\equiv X^{p^2}\pmod{(p,X^{p^2+1})}.
\]
It has the integral action of the specified quadratic integers
by \eqref{eq:integral-log}. Lift an isomorphism of the two
reductions by the height-one deformation theorem. Comparing
logarithms gives an integral isomorphism \(t(X)=l^{-1}(cL(X))\)
with \(c\in R^\times\) its tangent map. Nonzero prime-to-\(p\)
torsion on the good curve has integral affine coordinates in
\(R\): the corresponding torsion group scheme is finite
\'etale over the strictly henselian ring, so its distinct
sections reduce distinctly, including the zero section.
\cite{drinfeld-formal,serre-tate}

\end{proof}

Normalize the analytic invariants by the good-model homothety.
For the quasi-period homomorphism of the Weierstrass zeta
function write \(\eta(\omega)=s_2\omega+B\bar\omega\) on
\(\Gamma\).

\begin{proposition}[The integral norm relation]\label{prop:good-cm-norm}
For the normalized good model the quadratic correction \(s_2\)
belongs to \(R\). For every quadratic unit \(a\) at \(p\), its action in the good formal coordinate satisfies the congruence \eqref{eq:cm-norm-relation}
\[
 q([a]X)\sim a\sigma(a)q(X).
\]

\end{proposition}

\begin{proof}
\proofstep{Compare the isogeny divisor and its leading coefficient.}

The Legendre relation gives
\(B=\pi/\operatorname{area}(\C/\Gamma)\).
Then \(\varsigma_*(z)=\varsigma(z)e^{-s_2z^2/2}\) obeys
\[
 \varsigma_*(z+\omega)=\chi(\omega)
         \exp\big(B(z\bar\omega+|\omega|^2/2)\big)
         \varsigma_*(z),\qquad \chi(\omega)\in\{1,-1\}.
\]
For \(0\ne a\in{\cal O}_K\) let \(n=a\bar a\). On the complex
elliptic curve the following identity holds:
\begin{equation}\label{eq:cm-isogeny}
 \frac{\varsigma_*(az)^2}{\varsigma_*(z)^{2n}}
       =a^2\prod_{P\in\ker[a]\setminus\{0\}}
                                      (\wp(z)-\wp(P)).
\end{equation}
The left side is elliptic since the Hermitian exponent scales
by \(a\bar a\) and the signs square to one. The two sides have
the same divisor.

For a point with \(P=-P\) the factor on the
right has a double zero there; otherwise the two factors for
\(P,-P\) provide the double zero. Comparison of the first
Laurent coefficients at zero gives \(a^2\). Comparison at
the next even power gives
\[
                 (a^2-n)s_2=\sum_{0\ne P\in\ker[a]}\wp(P).
\]
Take \(a=\sqrt{-D}\). All these points have order dividing
\(D\), and \(a^2-n=-2D\) is a unit at \(p\). The good torsion
calculation therefore puts \(s_2\) in \(R\).

\proofstep{Use prime-to-\(p\) torsion and approximate the local units.}

For \(a\) a unit at \(p\), express \eqref{eq:cm-isogeny} in
the integral parameter \(t\). With \(A(t)=t^2\wp(l(t))\)
the result is
\[
 \exp\{-2(q_E([a]t)-nq_E(t))\}
  =\left(\frac{a t}{[a]t}\right)^2
                   \prod_{0\ne P\in\ker[a]}(A(t)-\wp(P)t^2).
\]
The right side is a principal integral unit. The factor
\([a]t/(a t)\) is integral with constant one by the extension
of the CM action, and all kernel points have prime-to-\(p\)
order.

Taking a square root is integral since \(p\) is odd.
Conjugate by \(t(X)\) and absorb the unit \(t(X)/(cX)\) as
in the preceding subsection. The norm relation
\eqref{eq:cm-norm-relation} follows for \(a\in{\cal O}_K\)
of prime-to-\(p\) norm. Such elements are dense in
\({\cal O}_{K_p}^\times\): approximate the two coordinates
in the integral basis \(1,(1+\sqrt{-D})/2\) modulo \(p^j\).
The first approximation modulo \(p\) already ensures the
norm is prime to \(p\).
Indeed, for the approximants \(a_j\) to a unit \(u\),
\[
  a_j\equiv u\pmod {p^j}\quad\Longrightarrow\quad a_j\bar a_j\equiv u\bar u\not\equiv0\pmod p .
\]
 On this quadratic field \(\sigma\)
is conjugation. Each coefficient of the exponential of the
norm difference is continuous in \(a,\sigma(a)\); its value
on all units is integral by closedness. This supplies the
comparison with all scalars in the trace calculation.

\end{proof}

Use the Heisenberg translation
\[
 (U_P f)(z)=f(z+P)\exp\{-B(z\bar P+|P|^2/2)\}.
\]
\begin{proposition}[Integral gauged torsion translates]\label{prop:good-cm-torsion-translates}
For each prime-to-\(p\) translate \(P\) of the normalized even
section \(\vartheta_0\), write \(G_P(t)=e^{q_E(t)}(U_P\vartheta_0)(l(t))\). The integrality and primitivity conclusions are
\[
 G_P(t)\in R[\![t]\!],\qquad G_P(t)\bmod\mathfrak m_R\ne0.
\]
 The same assertion holds in the
integral coordinate \(X\) with its gauge \(q\).
\end{proposition}

\begin{proof}
\proofstep{Obtain the addition recurrence.}

In \(t\) coordinates write \(G_P(t)=e^{q_E(t)}
 (U_P\varsigma_*)(l(t))\). In particular \(G_0=t\).
For \(P\) a period this is \(t\) times a sign. If \(Q\)
represents a nonzero torsion point, \(C=G_Q(0)\ne0\).
Sigma addition, with all quadratic exponential factors
cancelled, gives for \(P\) nonzero modulo periods
\begin{equation}\label{eq:torsion-difference}
      G_{P+Q}G_{P-Q}
            = C^2\{\wp(Q)-\wp(l(t)+P)\}G_P^2 .
\end{equation}
\proofstep{Determine the exact special-fibre orders.}

Take \(N\) the exact order of \(Q\), prime to \(p\). For
\(1\le j<N\) write \(f_j=\wp(Q)-\wp(l(t)+jQ)\).
It is an integral unit times \(t^{b_j}\), with
\(b_j={\bf1}_{j=1}+{\bf1}_{j=N-1}\). To check this, translation
by the integral section \(jQ\) identifies the completed local
ring there with \(R[[t]]\); \(x\) is regular on its neighbourhood.
The divisor \(x-x(Q)\) on the smooth Weierstrass curve has
zeros exactly at the sections \(Q,-Q\), counted with
multiplicity. On the special fibre its zeros have the same
description. Reduction injects this prime-to-\(p\) torsion
group; at order two the double zero is exact since the cubic
has distinct roots and \(p>2\). Thus the first nonzero
coefficient at the indicated order is a unit.

\proofstep{Recover the scalar and the integral series.}

Repeatedly using \eqref{eq:torsion-difference} gives
\[
 G_{NQ} G_0^{\,N-1}
       =G_Q^{\,N}\prod_{j=1}^{N-1}(C^2 f_j)^{N-j}.
\]
This follows by taking successive quotients \(G_{(j+1)Q}/G_{jQ}\)
in the Laurent series field. The total vanishing order of the
product of \(f_j\)'s is \(N\), including \(N=2\). The left side
is a unit scalar times \(t^N\). Cancel powers of \(t\) and
compare constants: 
\[
  |C|_p^{N^2}=1\quad\Longrightarrow\quad v_p(C)=0.
\]
 Hence \(G_Q^N\) is
an integral unit, and \(G_Q/C\) itself is one by the binomial
series with exponent \(1/N\). For \(Q\) representing zero the
formula \(G_Q=\pm t\) gives integrality directly.

Fix a
nonzero half-period \(Q_0\). Normalize the even translate
\(\vartheta_0=U_{Q_0}\varsigma_*/(U_{Q_0}\varsigma_*)(0)\).
Composing translations by \(P,Q_0\) differs from translation
by \(P+Q_0\) by a root of unity for torsion arguments, since
the alternating lattice pairing takes rational values on
division points. The denominator at \(Q_0\) is a unit by the
order-two case. Thus for any prime-to-\(p\) translate \(P\),
\(e^{q_E(t)}(U_P\vartheta_0)(l(t))\) is integral with some
unit coefficient. Substitution of the integral isomorphism
\(t(X)\) and the change \(q_E\mapsto q\) preserve this
assertion.

\end{proof}

\begin{auxlemma}[Extracting a coefficient or a derivative from disjoint copies]
\label{lem:cm-disjoint-tensor}
Let \(F\) be a valued field of characteristic zero and
\(H(X)=\sum_\alpha a_\alpha X^\alpha\in F[[X_1,\ldots,X_n]]\).
For any \(m\ge1\), take disjoint sets of variables \(X^{(1)},\ldots,X^{(m)}\)
and set \(T=\prod_{j=1}^m H(X^{(j)})\). For a multi-index \(\nu\),
\[
 \left[\prod_j(X^{(j)})^\nu\right]T=a_\nu^m,\qquad
 \left.\left(\prod_j\partial_{X^{(j)}}^\nu\right)T\right|_0
         =\bigl((\partial_X^\nu H)(0)\bigr)^m .
\]
Integrality of the indicated coefficient or derivative of \(T\)
therefore gives integrality of the corresponding coefficient or derivative of \(H\).
\end{auxlemma}

\begin{proof}
In the indicated monomial each set of variables has its own exponent \(\nu\),
so precisely one product of coefficients occurs. Multiplying each factor
by \(\nu!=\prod_i\nu_i!\) gives the derivative identity. If a nonzero
element has integral \(m\)-th power, then \(m\) times its valuation is
nonnegative. The zero reading is integral as well.
\end{proof}

\begin{proposition}[Integral correction in an inert CM chart]\label{prop:inert-cm-chart}
For the full-order curve at an inert prime \(p\ge5\),
use the good formal model after homothety. Suppose that, after an integral unitary change of orthogonal frame,
the fourth tensor power of the translated theta section is a finite
integral linear combination of products of the one-variable
translates of Proposition~\ref{prop:good-cm-torsion-translates}.
For formal coordinates \(\mathbf X\) on that chart, the corrected series satisfies
\[
 e^{\sum q}\Theta\in R[\![\mathbf X]\!].
\]
This includes \(\Theta(0)=0\). Here \(q\) is the quadratic correction
series from Section~\ref{subsec:curve-exponents}.
\end{proposition}

\begin{proof}
\proofstep{Transport the gauge by the integral unitary frame.}

To see where these one-variable units enter the orthogonal
chart, let \(\gamma\) be an integral unitary change of frame
over the quadratic integers at \(p\). Its inverse conjugate
transpose is integral. On \(4n\) coordinates it acts by
\(Z\mapsto L^{-1}(\gamma L(Z))\) componentwise. By the CM
norm relation and additivity of \(b\), the sum of \(q\) over
all coordinates is invariant modulo \(\sim\): the diagonal
contractions of \(\gamma^*\gamma\) are one and the off-diagonal
ones vanish.

For the contracted coefficients the two cases read
\[
  \sum_i\overline{\gamma_{ij}}\gamma_{ik}=(\gamma^*\gamma)_{jk}=\delta_{jk}.
\]
 Suppose in this frame the fourth tensor of the
translated theta section has a finite expansion in products
of the one-variable translates above, with integral coefficients.
After multiplication by the exponential of the summed \(q\)
each product is integral. Undo the unitary substitution and
absorb the principal unit from the invariant sum. The result
is that the fourth tensor of \(e^{\sum q}\Theta\) is integral.
\proofstep{Recover each coefficient from its fourth power in disjoint variables.}

Lemma~\ref{lem:cm-disjoint-tensor}, with \(m=4\) and 
\[
  H=e^{\sum q}\Theta,
\]
now gives integrality coefficient by coefficient, including at a zero
of the section.

The global scalar and its local standard frames are constructed in Section~\ref{subsec:theta-normalization}.
\label{cor:cm-models-1}
\end{proof}

\subsection{One normalization for the fourth theta tensor}
\label{subsec:theta-normalization}

Local integral models can have different generators. The fourth tensor needs a scalar normalization compatible with all the translates occurring in one period.

Fix an imaginary quadratic field \(K=\mathbb Q(\sqrt{-d})\)
with \(d\) prime, \(d\equiv7\bmod8\). Put
\(\tau=(b+\sqrt{-d})/2\) and take
\(\Gamma=\lambda(a\mathbb Z+\tau\mathbb Z)\), where
\(4ac=b^2+d\), \(a>0\) is odd and \(\lambda\ne0\).
Thus \(b\) is odd and \(c\) is even. Write
\(H_0(z,w)=\sum z_i\overline{w_i}/\operatorname{area}(\Gamma)\).

For a positive Hermitian form \(H\), linear in the first
variable, put \(E=\Im H\). A CM lattice \(M\) here is a full
\(\mathcal O_K\)-lattice self-dual for \(E\). Its sign character is
\[
  \alpha(x)=\exp\big(\pi i\,\sqrt d\,H(x,x)/2\big).
\]
Indeed the quadratic expression in the exponent divided by
\(\pi i\) is \(E(\tau x,x)\). For \(x,y\in M\) its
polarization is \(2E(\tau x,y)-bE(x,y)\), congruent to
\(E(x,y)\bmod2\). Consequently

\[
  \alpha(x+y)=\alpha(x)\alpha(y)e^{\pi i E(x,y)}.
\]
On one copy of \(\Gamma\) the quadratic expression is
\(au^2+buv+cv^2\) at \(\lambda(au+\tau v)\). Its parity is
\(u+uv\), the sign of the odd sigma section translated by
the half period \(\lambda a/2\). In an orthogonal sum the
signs multiply. Let \(f_\nu\) denote the translated one-variable sigma section in the \(\nu\)-th coordinate, divided by its nonzero value at zero. The product section in these coordinates is
\[
  F_0((z_\nu)_\nu)=\prod_\nu f_\nu(z_\nu),\qquad f_\nu(0)=1,\qquad F_0(0)=1.
\]

Take a positive CM Hermitian space of rank \(n\) with such
a lattice \(M\), in the complex scale of \(\Gamma\).
There is a rational isometry \(g_0\) of its fourfold sum to
the standard space of rank \(m=4n\). In fact diagonalize the
Hermitian form over \(K\) relative to the rational CM form
underlying \(H_0\). Every diagonal entry is a positive rational
number \(r\).

For a positive entry \(r=a/b\) with \(a,b\in\mathbb Z_{>0}\), Lagrange's theorem gives
\[
  ab=n_1^2+\cdots+n_4^2,\qquad r=\sum_{i=1}^4(n_i/b)^2.
\]
 The matrix of multiplication by the resulting
quaternion has transpose times itself equal to \(r I_4\).
Taking the block diagonal sum of the four-dimensional
quaternion matrices gives \(g_0\). Its equality of forms is
\[
       g_0^*H_0^{\oplus 4n}g_0=H^{\oplus4}
\]
in the corresponding ordered coordinates, with \(H_0\)
the standard form on one line. In particular \(L_0=g_0 M^4\) and \(S=\Gamma^m\)
are commensurable in the common rational space.

\paragraph{Choosing a local isometry and its determinant.}

To transport the normalized theta tensor from the standard lattice, we need a local isometry with the prescribed determinant.

\begin{lemma}[Local comparison with the standard CM lattice]
\label{lem:cm-standard-local-isometry}
For each finite prime \(p\) the lattices \(L_0\) and \(S\) are
locally isometric. The isometry at \(p\) can be chosen rational
with \(1-\delta\) invertible. For \(L_1=\delta L_0\), its intersection \(I=L_1\cap S\) satisfies
\[
 p\nmid[L_1:I]\,[S:I].
\]

\end{lemma}

\begin{proof}
After removing the common complex scale, write
\(\mathfrak a=a\mathbb Z+\tau\mathbb Z\) for the ideal of the
reference curve. It is invertible with

\[
  \mathfrak a\overline{\mathfrak a}=a\mathcal O_K,\qquad {\rm N}\mathfrak a=a
\]
 since
\(\mathrm N\mathfrak a=a\).

Tracing the Hermitian pairing
\(h(x,y)=\sum x_i\bar y_i/a\) against
\(1/\sqrt{-d}\) gives \(E(x,y)\). Since the inverse
different is \((1/\sqrt{-d})\mathcal O_K\), stability under
\(\mathcal O_K\) identifies self-duality for \(E\) with
unimodularity for \(h\). Locally choose a generator of
\(\mathfrak a\), so \(S\) has Gram matrix a unit times the
identity. At a split prime the two direct summands of a lattice
are in perfect duality. For a basis \(e_i\) of one summand and the dual basis \(e'_j\) of the other, the split pairing has entries
\[
  h(e_i,e'_j)=\delta_{ij}.
\]
This description includes \(p=2\).

At an inert odd prime, reducing a
unimodular form gives a nondegenerate Hermitian form over
\(\mathbb F_{p^2}\). It has a vector of nonzero norm by
polarization. Lift the vector, split its orthogonal complement
and iterate. The norm on unramified units is surjective, by
finite-field norm surjectivity and lifting the norm using the
surjective residue trace.
For the successive corrections to a norm, the linear term is the residue trace: for \(n\ge1\),
\[
  N(1+p^nx)\equiv 1+p^n\operatorname{Tr}(x)\pmod {p^{n+1}}.
\]
 Thus all diagonal units may be matched.

At the odd ramified prime the residual forms are nondegenerate
symmetric forms. They have the same determinant square class
because the generic change of basis has determinant of norm a
unit, and the residue of that norm is a square. Dimension and
determinant classify nonsingular forms over the finite field of
odd order. Lift a residual isometry modulo the uniformizer.
For a resulting Gram matrix \(B\) close to the desired \(C\),
\(B^{-1}C\) is self-adjoint for \(B\) and congruent to one.
Put \(A=B^{-1}C\) and \(R=A^{1/2}\).
The binomial series is integral since \(A\) is congruent
to \(1\) and \(1/2\in\mathbb Z_p\). Each power is
self-adjoint for \(B\), hence so is \(R\), and
\[
                  R^* B R=BR^2=BA=C .
\] This proves the local assertion.

The transporter of the two lattices is open in the local
unitary group. It contains elements for which both \(1+\delta\)
and \(1-\delta\) are invertible. For example multiply a
transporter by nearby scalar norm-one elements, avoiding its
finitely many forbidden eigenvalues. On this open subset
the Cayley change \(X=(1-\delta)(1+\delta)^{-1}\) parametrizes
the group by skew-Hermitian matrices. The rational skew-Hermitian
matrices are dense in their localization since they form a
rational vector space. We can therefore take \(\delta\) rational,
with \(\delta L_{0,p}=S_p\) and \(1-\delta\) invertible.
Put \(g_1=\delta g_0\). Each intersection index of
\(L_1=g_1M^4\) and \(S\) is prime to \(p\).
\end{proof}

For the choices in Lemma~\ref{lem:cm-standard-local-isometry}
one algebraic number suffices for normalization. Let
\(J=(\bigwedge^m L_0)/(\bigwedge^m S)\) be the fractional
ideal ratio in the common determinant line, with complex scale
removed. At every prime the local transporter has norm-one
determinant and \((\det\delta)J_p=\mathcal O_{K,p}\).
Hence \(J\bar J=\mathcal O_K\). Choose \(h>0\) with
\(J^h=(u)\), using finiteness of the ideal class group.
The positive rational number \(u\bar u\) generates the unit
ideal, so equals one. Fix any \(\kappa^{2h}=u\). All complex
conjugates of \(\kappa\) have modulus one. For every embedding
at \(p\) and every such \(\delta\) we have
\begin{equation}\label{eq:kappa-determinant}
             |\kappa|_p^2 |\det\delta|_p=1 .
\end{equation}
In particular the two valuations above a split prime are
accounted for separately.

\begin{lemma}[Integral Cayley change of the theta isometry]\label{lem:theta-cayley-scalar}
For the self-dual CM lattices with the intrinsic sign above,
use an isometry \(\delta\) for which the stated Cayley
operators are defined. The transport and its scalar satisfy
\[
 F\circ\delta^{-1}=c_\delta T_{L,\delta L}F,\qquad c_\delta^8=(\det\delta)^4.
\]

\end{lemma}

\begin{proof}
The theta line of a self-dual lattice with the given
semicharacter consists of the entire functions fixed by
\(\alpha(x)U_x\) for every lattice vector, where
\[
 U_xF(z)=e^{-\pi H(z,x)-\pi H(x,x)/2}F(z+x).
\]
It has dimension one by Riemann--Roch for a principal
polarization: the positive line bundle has vanishing higher
cohomology and top self-intersection divided by \(m!\) equal
to the Pfaffian of \(E\), of absolute value one. Equivalently
one can choose a symplectic basis and solve the automorphy
relations for its Fourier coefficients. For commensurable
self-dual \(L_i,L_j\) with the same quadratic sign on their
intersection, define
\[
 T_{ij}F=[L_j:L_i\cap L_j]^{-1/2}
       \sum_{x\in L_j/(L_i\cap L_j)}\alpha(x)U_xF .
\]
The Heisenberg multiplication rule

\[
  U_x U_y=e^{-\pi i E(x,y)}U_{x+y}
\]
 shows independence of
representatives and the fixed-vector property on \(L_j\).
Use the Gaussian metric and normalized Haar average on any
common period domain. The distinct summands are orthogonal:
translation by \(\ell\in L_i\) distinguishes the characters
\(e^{2\pi i E(x,\ell)}\), by self-duality of \(L_i\).
Each has the norm of \(F\), since the translated density has
the same mean. Thus \(T_{ij}\) is an isometry of lines.

For three such lattices carrying the intrinsic quadratic signs,
\(T_{23}T_{12}=\xi T_{13}\) with \(\xi^8=1\). To compute the
phase, project onto the initial line by averaging its
translations, or equivalently retain the coefficient of the
trivial \(L_1\)-character in the preceding orthogonal expansion.
Up to positive index factors the scalar is the sum over pairs
\(z\in L_3,y\in L_2,z+y\in L_1\) of
\[
 \alpha(z)\alpha(y)\alpha(z+y)e^{-\pi i E(z,y)} .
\]
For a common finite quotient take \(z,y\) modulo
\(2(L_1\cap L_2\cap L_3)\) in their respective lattices.
This counts each contributing pair of cosets in the two
intertwiners with the same multiplicity: translates of
\(z\) by \(t\in L_2\cap L_3\) are paired with translates of
\(y\) by \(-t\), and \(y\) can independently be shifted by
\(L_1\cap L_2\).

The Heisenberg law leaves the summand
unchanged under both operations. The displayed function
on the finite group of pairs is a homogeneous quadratic
character, by its quadratic exponent. Its sum is nonzero
because the intertwining maps are isometries and the
coefficient of the initial line in \(T_{13}\) is positive.
On the radical of its polarization its restriction is a
character.
Indeed on the radical \(R\) the defining polarization is trivial, so
\[
  q(x+r)=q(x)q(r)\qquad(r\in R).
\]
 A nontrivial such restriction would annihilate
the sum by translation.
In fact translation by a radical element \(r\) multiplies the sum by \(q(r)\), giving
\[
  \sum_xq(x)=\sum_x q(x+r)=q(r)\sum_xq(x).
\]
 Passing to the quotient by the
radical therefore changes the sum by a positive integer.

For completeness, let \(q\) be a nondegenerate homogeneous
quadratic character on a finite abelian group \(B\). Writing \(b(x,y)=q(x+y)/(q(x)q(y))\), character orthogonality gives
\[
  \left|\sum_xq(x)\right|^2=\sum_h q(h)\sum_y b(h,y)=|B|.
\]
Denote this Gauss sum by \(G\). To see that its
phase is an eighth root of unity, put \(N=|B|\) and choose
four integer squares summing to \(N^2-1\). Left quaternion
multiplication by this vector gives an integer similarity
\(Q\) of the sum of four squares with
\(Q^{t}Q=(N^2-1)I\), invertible on \(B^4\).
Expanding the quadratic character and its symmetric
bicharacter gives

\[
  \prod_i q((Qv)_i)=(\prod_i q(v_i))^{N^2-1}
=\overline{\prod_i q(v_i)}.
\]
 Here \(q(Nx)=q(x)^{N^2}=1\).
Thus \(G^4\) is real. This also covers \(N=1\).
As \(|\xi|=1\), the triple comparison follows.

An additional determinant occurs for a change of variables.
Let \(L=S\) and take the rational unitary \(\delta\) of
Lemma~\ref{lem:cm-standard-local-isometry}. Put
\(A=1-\delta^{-1}\), \(I=L\cap\delta L\) and \(j=[L:I]\).
Then for \(F\) on the theta line of \(L\),
\begin{equation}\label{eq:cayley-multiplier}
   F\circ\delta^{-1}=c_\delta T_{L,\delta L}F,\qquad
              c_\delta^8=(\det\delta)^4 .
\end{equation}

To compute the scalar, periodize the Gaussian kernel in its
first variable:
\[
 P(v,z)=\sum_{x\in L}\alpha(x)
      U_x\big(e^{\pi H(\cdot,z)}\big)(v).
\]
Positive definiteness gives normal convergence on compact sets
by a Euclidean lattice Gaussian, and \(P(\cdot,z)\) lies in the
theta line. Integrate \eqref{eq:cayley-multiplier}'s first
identity for this section, evaluated at \(z\), against
\(e^{-\pi H(z,z)}\,dz\) on a fundamental cell of \(I\).
Before inserting \(c_\delta\), the right side integrates to
\(\operatorname{vol}(I)/\sqrt j\), where \(\operatorname{vol}\)
of a lattice denotes covolume. Indeed in the expansion of
\(T P\) the oscillatory factors are characters
\(e^{2\pi i E(x+y,z)}\) for \(x\in L\), \(y\in\delta L/I\);
integration leaves just \(x+y=0\).

On the left the summand for \(x\) is
\[
 q_A(x)\,e^{-\pi H(A(z-A^{-1}x),\,z-A^{-1}x)},\qquad
 q_A(x)=\alpha(x)\exp[-\pi H((A^{-1}-\tfrac12)x,x)] .
\]
This is completion of the square using
\(A+A^*=A^*A\). Its absolute Gaussian has exponent
\(-\pi H(Az,Az)/2\) after recentering. The polarization
of \(q_A\) on \(L\) is \(e^{-2\pi i E(A^{-1}x,y)}\).
Its radical consists precisely of \(AI\): self-duality
says \(A^{-1}x\in L\), and \(x\in L\) forces
\(\delta^{-1}A^{-1}x\in L\). For \(x=Av\) with \(v\in I\),
substituting \(x=v-\delta^{-1}v\) and

\[
  \alpha(\delta^{-1}v)=\alpha(v)
\]
 into the semicharacter
law gives \(q_A(x)=1\). Hence it descends to a
nondegenerate homogeneous quadratic character on \(L/AI\).
Group the sum by this quotient and unfold the translates by
\(I\). Absolute integrability follows from the finite number
of recentered Gaussians. Diagonalizing \(A\) in unitary
coordinates yields
\[
 \int e^{-\pi H(Az,z)}\,dz=
      (\det A)^{-1}\int e^{-\pi H(z,z)}\,dz
       =\operatorname{vol}(L)/\det A .
\]
Each eigenvalue has positive real part since \(1\) is excluded
from the spectrum of \(\delta\). The last equality follows in
orthonormal coordinates from the symplectic covolume one of a
self-dual lattice. Since \([L:AI]=j|\det A|^2\), division by
the right-hand integral gives
\[
  c_\delta=\frac{|\det A|}{\det A}
      \frac{1}{\sqrt{[L:AI]}}\sum_{x\in L/AI}q_A(x).
\]
The normalized sum has eighth power one.
Finally \(A^*=-\delta A\), so
\(\overline{\det A}=(-1)^m\det\delta\det A\);
this proves the asserted eighth power with its sign and exponent.
\end{proof}

\paragraph{Returning to the local coefficients.}
Apply Lemma~\ref{lem:theta-cayley-scalar} before taking the
fourth tensor and returning to the local argument lattice.
Set \(K_i=(T_{S,L_i}F_0)\circ g_i\). Transport of \(U_x\)
under \(\delta\) and the triple comparison give
\(K_0=cK_1\) with \(c^8=(\det\delta)^4\). Explicitly, in
the \(g_1\) coordinates the transported section is
\(T_{\delta S,L_1}(F_0\circ\delta^{-1})\).
Use \eqref{eq:cayley-multiplier} and then compose the two
intertwiners. Since the theta line of \(M^4\) is the fourth
tensor power of the line of \(M\), there exists a theta
vector \(\Theta\), unique up to a fourth root of unity,
with \(\Theta^{\otimes4}=\kappa K_0\). This fixes its
complex norm at every conjugate of the determinant scalar.

For each finite embedding 
\[
  \Theta^{\otimes4}=\kappa K_0=(\kappa c)K_1
\]
and \(|\kappa c|_p=1\) by \eqref{eq:kappa-determinant}.
Every summand of \(T_{S,L_1}F_0\) is a product of
one-variable translates with prime-to-\(p\) denominators,
and its normalization is a sign times a unit square root
of an inverse index. A further translation by a point of
order prime to \(p\) adds points of the same kind in the
standard coordinates: a prime-to-\(p\) multiple of its
\(g_1\)-image lies in \(L_1\), and a prime-to-\(p\) index
carries that lattice into \(S\).

Since the two translation arguments belong to the rational space of \(S\), write \(E(y,y')=a/b\) with \(b>0\). The composition phase satisfies
\[
  \bigl(e^{\pi iE(y,y')}\bigr)^{2b}=e^{2\pi ia}=1.
\]
This gives the normalized orthogonal fourth tensor expansion
used in Section~\ref{subsec:cm-models}. Section~\ref{subsec:baseline-jets} expresses the
coordinate change on the argument lattices by CM matrices
integral at the place and deduces integrality of derivatives.
The positive gains also use the incident chart and its
linear forms in Section~\ref{subsec:incident-sections}.

\subsubsection{Integral derivatives away from the discriminant}
\label{subsec:baseline-jets}
\hypertarget{source-baseline-jets}{}

Beyond the primes where the jets give a positive gain, the product of local estimates still requires integrality in every order of differentiation. This applies at the level-free places. We keep the
scale of the curve ideal used in the period. Thus the period
lattice \(\Gamma\) has discriminant
\(g_2(\Gamma)^3-27g_3(\Gamma)^2=1\). Use the basis
\(\lambda a,\lambda(b+\sqrt{-d})/2\) of
Section~\ref{subsec:theta-normalization}, with \(a\) odd.
Write \(E=\C/\Gamma\), \(h=\lambda a/2\) modulo \(\Gamma\),
and \(\theta\) for its even degree-one section normalized by
\(\theta(0)=1\). Its sole zero is \(h\).
Fix a coefficient embedding into \(\mathbb C_p\) and extend it
to a field isomorphism on \(\C\) as in
Section~\ref{subsec:word-measure}. Elliptic functions will
be transported as rational functions on the Weierstrass curve.
All formal identities use the transport on coefficients.

\begin{proposition}[Integral jets at level-free primes]\label{prop:level-free-theta-jets}
Use a discriminant-one curve with period lattice \(\Gamma\),
odd \(a\) in the basis of Section~\ref{subsec:theta-normalization},
and the degree-one even section normalized by \(\theta(0)=1\)
with zero at \(\lambda a/2\bmod\Gamma\).
At \(p\ne d\), on the admissible lattice pairs with the
stated level exclusions and prime-to-\(p\) denominators,
let \(a_{\alpha,w}\) be the tested ordinary derivative of multi-order \(\alpha\) at a coefficient place \(w\). Then
\[
 v_w(a_{\alpha,w})\ge0\qquad\text{for every permitted }\alpha,w.
\]
 The bound is uniform in the order of
the derivative and in the coefficient field. It includes
the primes \(2\) and \(3\) when they split in the chosen CM field.
\end{proposition}

\begin{proof}
\proofstep{Choose a good model with the given invariant differential.}
The CM \(j\)-invariant is an algebraic integer. Consequently
the curve has potentially good reduction. To keep track of
the differential, take a good model after extension and choose
its analytic homothety \(c\) from the fixed lattice. Its
discriminant is \(c^{-12}\), so \(c\) has norm one. The change
of Weierstrass variables with zero translations and unit scaling
returns an integral good model with the original \(g_2,g_3\)
invariants. Over \(\mathbb C_p\) the normalized differential
therefore identifies the formal logarithmic variable with the
transported analytic variable \(z\). Write \(T\) for a good
formal parameter and \(z=l(T)\) for the formal logarithm,
with unit tangent. In particular

\[
  D_z=(l'(T))^{-1}\frac{d}{dT},\qquad D_z\big(R[\![T]\!]\big)\subset R[\![T]\!]
\]
 preserves integral formal series
because the invariant differential has integral unit coefficient
in this parameter. For existence over a finite local extension,
the invariants \(g_2,g_3\) are algebraic by the discriminant
identity and \(j=1728g_2^3\). The analytic scale fixes the
ratios of the invariant differentials. \cite{cm-reciprocity,serre-tate}

\proofstep{Check that the odd characteristic remains visible in the special fiber at two.}

For odd \(p\), \(h\) has nonzero reduction by injectivity on
prime-to-\(p\) torsion. At two we use its particular
characteristic. The ring \(O=\Z[(1+\sqrt{-d})/2]\) splits at two,
since \(d\equiv7\pmod8\). In \(O/2O\) the integers
\(\omega=(1+\sqrt{-d})/2\) and \(\bar\omega\) give the
two complementary idempotents. The two idempotents split the point as
\[
       h=\omega h+\bar\omega h,\qquad
       \omega h\ne0,\quad\bar\omega h\ne0\quad\text{in }E[2].
\]
Indeed the second lattice coordinate of the corresponding lift
is \(a/2\) or \(-a/2\), which excludes membership in \(\Gamma\).

Each idempotent cuts out a height-one summand of
\(E[2^\infty]\), as one sees on the generic Tate module.
The action extends over a good model. The special fiber has
dimension one, so one summand has dimension zero and is
\'{e}tale, and the other has dimension one and is connected.
Reduction is injective on points of the \'etale summand.
Thus the reduction of \(h\) has order two at this place as
well.  It holds for either embedding of \(O\).

\proofstep{Use a split dilation with a unit tangent and norm divisible by the residue prime.}

Suppose that \(p\) splits in \(O\), including \(p=2\) and
\(p=3\) in the selected host (\(d\equiv2\pmod3\)).
Order the components of \(O\otimes\Z_p\) with the differential
embedding first.
At a split prime the two idempotents are
\[
  O\otimes\mathbb Z_p\simeq\mathbb Z_p\times\mathbb Z_p,\qquad e_1=(1,0),\quad e_2=(0,1).
\]
 Choose \(\alpha\in O\) congruent to \((1,0)\)
modulo \(p\), nonzero and of norm \(m>1\). Its action
\([\alpha](T)\) on the good formal curve is integral and
has unit linear coefficient \(A\), while \(p\mid m\).
The quadratic sign of the theta vector at a lattice period
\(v\) is \((-1)^{E_0(\omega v,v)}\), with \(E_0\) its alternating
form, as in Section~\ref{subsec:theta-normalization}.

Replacing
\(v\) by \(\alpha v\) multiplies that integer by \(m\).
The Hermitian automorphy exponents also multiply by \(m\).
The quotient
\[
             R(z)=\theta(\alpha z)/\theta(z)^m
\]
is therefore elliptic. Its divisor is supported on \(h\) and
the preimages of \(h\). None reduces to the identity since
\(\bar h\ne0\) and \([\alpha]\) preserves reduction. Both
\(R\) and its inverse have unit integral expansions at the
identity, with constant term one. For detail, its divisor is
defined over a finite extension of the good reduction field,
so the rational function descends there up to a scalar.
A degree-zero principal divisor of rational points has a
one-dimensional space of defining rational functions, by base
change of sections on the proper curve.

If two such functions \(f,g\) have the prescribed divisor and are normalized at zero, then
\[
  {\rm div}(f/g)=0,\qquad (f/g)(0)=1\quad\Longrightarrow\quad f=g.
\]
 The normalization at
zero fixes the scalar. On a smooth model the only possible
component of the divisor through the origin is vertical.
After dividing by the corresponding power of a uniformizer
one has a unit in the local ring, since that ring is regular
and all divisor orders there vanish. Evaluation at the origin
shows that the original vertical order was zero.

Put \(F(T)=\theta(l(T))\) using the fixed good-model coefficient
embedding into \(\mathbb C_p\) and the normalized \(\theta(0)=1\). The identity
\[
             F([\alpha](T))=F(T)^m R(l(T))
\]
gives integrality of \(F\) from its constant term one. At degree
\(j>0\) the unknown coefficient occurs with multiplier
\(A^j-m\), which is a unit, and all remaining terms are
integral by induction. The same conclusion holds for the
inverse series. This argument uses the split CM action on the
ordinary curve, including in residue characteristic three
where changing to a short integral equation could lose a
factor of three. At three the CM norm identity of \eqref{eq:cm-isogeny} gives the same integral correction with
\(\alpha=\sqrt{-d}\). There the correction

\[
  (\alpha^2-1)b_2/12=-((d+1)/3)b_2/4
\]
 is integral for the
general good equation. \label{cor:baseline-jets-1}

\proofstep{Use the logarithmic derivative of the theta section for primes at least five.}

For \(p\ge5\) with \(p\ne d\), the short equation at this
scale is good. Indeed \(j\) integral gives \(g_2\) integral
and the discriminant identity then gives \(g_3\) integral.
Apply \eqref{eq:cm-isogeny} with \(\alpha=\sqrt{-d}\).
Its kernel is killed by \(d\), hence every nonzero kernel point
has integral affine coordinates. The quadratic-coefficient
identity following \eqref{eq:cm-isogeny} gives \(s_2\) integral
because \(\alpha^2-\alpha\bar\alpha=-2d\) is a unit. The
second logarithmic derivative of the even theta function is
\[
           (\log\theta)''(z)=-\wp(z+h)-s_2 .
\]
The shift by \(h\) and the scalar gauge in the translated sigma
function contribute only a linear term to the logarithm in
addition to the displayed quadratic correction. The translated
Weierstrass coordinate is regular integral on the identity
disc by good reduction and \(\bar h\ne0\). All of its invariant
derivatives at the origin are integral.

Since \(\theta(0)=1\)
and \(\theta'(0)=0\), every positive derivative of
\(\log\theta\) at zero is integral. Applying Leibniz's rule
successively to \(\theta'=\theta(\log\theta)'\) proves
integrality of every ordinary derivative of \(\theta\).
This also treats inert primes where the argument with the
split dilation is unavailable.

\proofstep{Remove the zero of a translated section before taking a prime-to-residue-characteristic root.}

Let \(t\) have order dividing \(M\) prime to \(p\). With the
normalized translation \(U_t\) above, put
\[
       Q_t(z)=\left(\frac{U_t\theta(z)}{\theta(z)}\right)^M .
\]
The ratio has a period character of order dividing \(M\) by
the Heisenberg commutator. Thus \(Q_t\) is rational on the
curve and its divisor is \(M([h-t]-[h])\).
Let \(|\cdot|_{\rm vert}\) be the norm at the
generic point of the good special fibre. The telescope
of theta ratios gives
\[
          \prod_{j=0}^{M-1} Q_t(z+jt)\ \in\ \mu_\infty:
\]
after cycling the translates the only factor is a
Heisenberg phase and the \(Mt\)-period sign. Each
translation preserves the special fibre and its generic
valuation. Hence \(|Q_t|_{\rm vert}^M=1\) and
\(|Q_t|_{\rm vert}=1\), including any scalar of the
normalized theta ratio by multiplicativity.

At \(p=2\) the reduction of \(h\) has order two and \(t\)
has odd order. At the odd primes the points generated by
\(h,t\) are prime-to-\(p\) torsion and reduction is injective
on them. Thus the divisor on the identity disc has either
no support or a zero of order \(M\) at the identity, the
latter precisely when \(t=h\) on the curve. The order on the identity disc can thus be written
\[
  \operatorname{ord}_{T=0}(Q_t)=eM,\qquad e=\begin{cases}1&t=h,\\0&t\ne h.\end{cases}
\]
 The expansion of \(Q_t/T^{eM}\)
is an integral unit with unit constant term. To see this, take
a defining function over a finite extension up to a scalar,
remove its vertical order, then remove the power of a regular
parameter defining the identity section.

The remaining
function and inverse are units at the special point by the
local divisor criterion above. The original generic norm
one makes the residual scalar a unit. The formal root
\((U_t\theta/\theta)(l(T))/T^e\) thus has unit constant and
integral coefficients. After dividing by its constant term,
write the normalized root as \(1+\sum_{j\ge1}b_j T^j\).
Then
\[
 [T^k]\Bigl(1+\sum_{j\ge1}b_j T^j\Bigr)^M
      =M b_k+P_k(b_1,\ldots,b_{k-1})
             \quad(P_k\in\mathbb Z[b_1,\ldots,b_{k-1}]).
\]
Since \(M\) is a unit, induction gives integral \(b_k\)
at every positive degree. The ratio itself is an integral series in \(T\).
Consequently its invariant derivatives are integral. The
product rule and the preceding bounds for \(\theta\) prove
the unit bound for all jets of \(U_t\theta\). In particular
the calculation covers the translate whose constant term
vanishes without dividing by that value.

\proofstep{Transfer the unit bound through both integral lattice comparisons.}

Consider a kernel term between the actual source and target
lattices. Its admissible pair agrees locally with these
lattices at \(p\). Choose positive integers \(s_1,s_2\)
prime to \(p\) clearing their two lattice comparisons in both
directions. For every test matrix \(y\) mapping the source
into the target, \(s_1s_2 y\) maps the admissible source into
the admissible target.
Neither clearing factor affects the valuation of a coordinate of a test direction, since
\[
  v_p(s_1s_2)=0,\qquad v_p(s_1s_2v)=v_p(v).
\]
 Multiplication by any period of
\(\Gamma\) sends this scaled embedded direction into the
period lattice of the Hom space. Likewise a translate
prime-to-\(p\) on the original pair has a prime-to-\(p\)
denominator on the admissible pair, by the same inclusions.

Apply the standardizing isometry \(g\) of
Section~\ref{subsec:theta-normalization} to the four copies.
It carries their lattice into the standard one after multiplication
by an integer prime to \(p\), by the relative index.
For each coordinate \(v\) of a scaled test direction the same
integer \(b_0\), prime to \(p\), gives
\[
 b_0v\Gamma\subseteq\Gamma
 \quad\Longrightarrow\quad b_0v\in\operatorname{End}(\Gamma)=O
 \quad\Longrightarrow\quad v\in O\otimes\Z_{(p)}.
\]
Here the endomorphisms are complex multipliers of the period lattice.
This proves integrality of the argument matrix under the chosen
embedding. Undoing the two lattice comparisons only multiplies
its columns by \(1/(s_1s_2)\). Further integral linear changes
of the two local bases, including the dual basis on the source,
preserve this conclusion by the chain rule.

In the four copies the normalized section is a finite sum of
products of the degree-one translates just treated, with
integral scalars. Integral ordinary derivatives are stable under
such sums, products and the indicated linear substitutions.
For instance the product step uses the integral coefficients in
\[
  \partial^\alpha(FG)(0)=\sum_{\beta\le\alpha}\binom{\alpha}{\beta}\partial^\beta F(0)\,\partial^{\alpha-\beta}G(0).
\]
The derivative identity in Lemma~\ref{lem:cm-disjoint-tensor}, applied
to the normalized section with \(m=4\), gives the unit bound for each
ordinary derivative in one copy. A prescribed
integral linear combination of translates satisfies that
bound as well. Use the stated admissible lattices and level exclusions and the prime-to-\(p\) denominator throughout. The estimate has constant one independently
of derivative order, the number of terms, the word and the
degree of the coefficient field. \label{cor:baseline-jets-2}
\end{proof}


\subsection{The period entries without a neighbor operator}
\label{subsec:unserved-period}
\hypertarget{source-unserved-period}{}

We first apply the jet and normalization estimates to entries for which no neighboring lattice is inserted. These orders provide the comparison for the entries modified by Hecke operators later in the chapter.

\begin{established}[Places and diagram sizes for the period estimate]\label{hyp:unserved-level-free}
At stage \(j\) use the valuation \(v=v_w/e_w\) for \(w\mid p\),
and let \(H_z\) be the number of boxes at the real label \(z\) with
positive stage weight \(w_z\). Set \(\kappa=k-1\), \(M=N_m\).
For the reference row and the row assigned to \(w\),
\[
 \begin{gathered}
 N_{j-1},N_j\le M,\quad 2\mid N_{j-1},N_j,\quad X''\le a,\quad
 \kappa w_z\le H_z\le a^2,\\
 0\le H_z-\kappa w_z\le4w_z\quad (z\text{ outside the tops}).
 \end{gathered}
\]
Call the place level free for \(\mu\) if
\(p\notin\{T_u,g_0,g_1\}\cup Q(\mu)\), with \(Q(\mu)\) the underlying
twist primes. The ordinary calculation here is restricted to
\(p\nmid N_{\rm host}g_0g_1\) and \(p\le X''\).
\end{established}
\begin{proof}
In \(N_i=k+i(t+1)\) the integer \(k\) is even and \(t\) is odd.
Also \(X''/a=nX'/b\to0\). The real labels are the \(f\) pairs
\((\pm,t)\) of the weight construction. Write \(e=w_z\) and \(N=N_m\). The head coordinates and rank shifts
are those calculated before the integral-kernel comparison in
Proposition~\ref{prop:level-four-integral-kernel}. Off the tops, if \(u\) ranges over the interior Harish--Chandra entries, then
\[
 u_{\max}=4(k-2)\Delta+4k\Delta+\tfrac12,\qquad
 u_{\min}=-\tfrac12,\qquad |u|\le4k\Delta .
\] Thus each of the \(k\)
head rows differs from \(e\) by at most \(S_c+2N\), where one may take
\(S_c=16k\Delta\). Each of the at most \(m(t+1)\) tail rows lies
between zero and \(2\eta\). With the common error
\[
 U=k(S_c+2N)+2m(t+1)\eta,\qquad |H_z-ke|\le U
          \quad(z\text{ off the tops}) .
\]

Here the signed stage offsets and appended rows use the same weight
coordinates as the explicit degree sum in Section~\ref{subsec:selected-exponent}.
On the selected range \(e\ge w''=b/(nX')\), while
\(U\le x^{O(1)}\exp(x^{\epsilon/4})\). Therefore
\[
 U/e\le x^{O(1)}
       \exp(x^{\epsilon/4}+x^{\epsilon/3}-x^\epsilon)\longrightarrow0 .
\]

This proves \((k-1)e\le H_z\le(k+1)e\) off the tops. At a top, delete
the exceptional first row. The \(k-2\) interior head rows have the same
bound as above and all the remaining rows are nonnegative. The deleted
row has at least \(a_t+e-N\) boxes. Since \(a_t\ge xb\) and \(e\le3b\),
for large \(x\)
\[
 H_z-\operatorname{row}_1\ge(k-2)e-U\ge(k-3)e,\qquad
 \operatorname{row}_1\ge2e .
\]

The weights are \(O(a)\) and the number of rows grows polynomially
in \(x\), so their total is at most \(a^2\) beyond one common threshold.
The assigned tuple avoids its coefficient prime when its twist places
are chosen; this is the prime avoidance in the construction of
Proposition~\ref{prop:translate-completion}. The two type primes are
distinct from the discriminant prime.
\end{proof}

\paragraph{The finite sum on a class pair.}
\begin{proposition}[Evaluation of the finite kernel on a pair of classes]\label{prop:unserved-reader-integral}
At a level-free place of Lemma~\ref{hyp:unserved-level-free}, the value of the finite kernel on this pair of classes is the sum
\eqref{eq:unserved-reader}. For each coefficient \(a\) and translation point \(P\) occurring there,
\[
 v_w(a)\ge0,\qquad \gcd(\operatorname{ord}P,p)=1,
\]
  and
the Taylor coefficients in the prescribed differential frames
can be transported coefficientwise by the embedding at \(w\).
\end{proposition}

\begin{proof}
\proofstep{Use the local levels of the two self-dual companions.}

Fix both classes of the matrix entry. Write \(L,L'\) for their
companion lattices in the two standard Hermitian spaces. They are
self-dual by the localization argument of
Proposition~\ref{prop:incident-companion-lattices}. Let \(\Lambda\) be the
embedded period lattice obtained from the inverse half different
times \(\operatorname{Hom}_{\mathcal O_E}(L,L')\), with the
curve period scale. At the finite test places the two half lattices
\(A,B\) have opposite level exponents.

If \(c\) is the positive
integer defining the level of the finite kernel, its prime divisors
are in \(S=\{T_u,g_0,g_1\}\cup Q(\mu)\), and
\[
 cA+cB\subseteq\Lambda,\qquad c\Lambda\subseteq A\cap B.
\]
These follow by comparing each of the exponents with the
valuation of \(c\) locally, including the zero exponents off the
level. Write \(S_X,S_A,S_B\) for representatives, respectively, of
\[
 A/(A\cap\Lambda),\quad
 (\Lambda\cap A)/(\Lambda\cap A\cap B),\quad
 B/(B\cap\Lambda).
\]
In particular \(h=|S_A|\) is prime to \(p\). Indeed the middle
quotient is finite and killed by \(c\).

\proofstep{Compute the finite Fourier expansion.}

Put \(E_H(x,y)=\Im H(x,y)\) for the alternating period form,
\(\alpha(x)=\exp(\pi i\sqrt d\, H(x,x)/2)\), and let
\(\Xi\) be the actual finite test extended by zero off the
embedded rational Hom space. For the normalized theta
\(\Theta_{L,L'}\), the finite part of the entry is
\begin{equation}\label{eq:unserved-reader}
 \sum_{x\in S_X,\ y\in S_B}
   \alpha(x)\alpha(y)e^{-\pi i E_H(x,y)}
   \left(\frac1h\sum_{s\in S_A}
      \Xi(s-x)e^{2\pi i E_H(s,y)}\right) U_{x+y}\Theta_{L,L'} .
\end{equation}
The Fourier sum at \(y\) uses the inverse commutator \(\exp(-2\pi i E_H(s,y))\), with the counting
measure divided by \(h\). The kernel reading uses volume one.
Applying \(U_x\) to the resulting \(U_y\) gives the displayed
Weyl factor. Since \(E_H(\Lambda,\Lambda)\subset\mathbb Z\)
and \(\sqrt d\,H(v,v)/2\in\mathbb Z\) on \(\Lambda\), all three
phase factors and the Fourier character are roots of unity.
The translates have order dividing \(c\) modulo \(\Lambda\).

\proofstep{Check the additive and determinant character factors.}

The word support is
a characteristic function. The type tests and twist tests on
a residue matrix multiply finite additive characters and values
of finite multiplicative characters, with zero on the excluded
residue classes. Each nonzero value is a root of unity.
Thus, for a nonzero finite character factor \(\xi\) of order \(n\),
\[
  n\,v_w(\xi)=v_w(\xi^n)=v_w(1)=0 .
\]
 In a
moved test the extra factors are integer powers of square roots
of the residue cardinality and local determinant characters.
At a type place the two movers preserve the submitted lattices,
so their determinants are local units and the residue characters
suffice.

At a twist place the uniformizer value and its inverse
are integral away from that place's rational prime and the type
levels. Indeed raise the ideal of the place to a principal power
with generator congruent to one at the character conductor.
The defining algebraic infinity type expresses that power of
the value as a product of conjugates of the generator to
integer powers, up to a root of unity. If \(\mathfrak P^A=(\alpha)\) with \(\alpha\equiv1\)
at the conductor, the value \(u\) satisfies

\[
  u^A=\zeta\prod_\tau\tau(\alpha)^{a_\tau}
\]
in the algebraic normalization. Thus
\[
       A v_w(u)=\sum_\tau a_\tau v_w(\tau(\alpha))=0
\]
for \(w\) outside that rational prime and the type levels. The off-type matching of the
column character with this ideal character gives the assertion
for the actual local value. The types are disjoint from the
auxiliary twist places. Finally the pinned terminal row either
uses an indicator or its finite Fourier reading. Its nonzero
values are integral additive character values, with a possible
factor \(q^{-Mn}\) at a place of cardinality \(q\), rank \(n\)
and depth \(M\). At the coefficient place under consideration, \(q\) is a unit. For any nonzero additive-character value \(\xi\) this gives
\[
  v_w(q^{-Mn}\xi)=-Mn\,v_w(q)+v_w(\xi)=0.
\]
 The
nonterminal classes use the matrix test, by the terminal pinning
of the supplementary rows. This proves integrality of every
scalar multiplying a translate in \eqref{eq:unserved-reader}.
It applies also when an entire entry vanishes.
\label{cor:unserved-kernel-charts-1}

All Taylor identities here are identities of coefficients.
\proofstep{Transport the Taylor coefficients in their fixed normalization.}

For clarity, algebraicity of the translated derivatives can
also be tested at one auxiliary good prime. Choose it away from
the denominators of the period comparison, the translate and a
rational basis of the Hom directions, away from the discriminant,
and split completely in the host. The fourfold expansion and the
baseline derivative estimate of Section~\ref{subsec:baseline-jets}
apply for every isomorphism of \(\mathbb C\) with \(\mathbb C_p\)
over a fixed embedding of \(\overline{\mathbb Q}\) at this auxiliary
prime. For each fixed derivative \(c\) in those directions the resulting bound \(M_c\) is uniform in the chosen isomorphism:
\[
  |\iota(c)|_p\le M_c\qquad\bigl(\iota:\mathbb C\xrightarrow{\sim}\mathbb C_p\text{ over the fixed algebraic embedding}\bigr).
\]
 In particular, the argument matrices are
integral after a fixed denominator clear for the rational
directions. The embedded matrices of a rational basis over the
quadratic field span the complexified Hom at the CM type, by
separability of \(E/K\).

For a fixed total order \(d\), the induced action on derivatives
is the degree-\(d\) tensor of the inverse direction matrix. Its
entries are algebraic and fixed. Thus every labelled coefficient
is a fixed finite algebraic linear combination of the bounded
directional derivatives. The triangle inequality gives the same
uniformity in the field isomorphisms.
Suppose one coefficient \(c\) were transcendental, and fix one
of the comparison isomorphisms \(\iota\colon\mathbb C\to\mathbb C_p\).
For each \(n\ge1\), the map 
\[
  c\mapsto p^{-n}c
\]
 extends by
transcendence bases to an automorphism \(\alpha_n\) over
\(\overline{\mathbb Q}\). The same directional bound would then
apply to all the isomorphisms \(\iota\circ\alpha_n\), whereas
\[
 |(\iota\circ\alpha_n)(c)|_p=p^n|\iota(c)|_p\longrightarrow\infty .
\] Division by the multi-index factorial preserves
algebraicity. The finite sums, linear substitutions and the
map on derivatives can consequently be transported by the
specified coefficient embedding at the original \(w\).

\end{proof}

\paragraph{Removing the differential generator in the actual frames.}

A differential generator can change the valuation in a scalar chart. In the actual integral frames the following comparison removes that factor.

\begin{proposition}[Removing the generator in integral frames]\label{prop:unserved-normalized-frame}
At the enhanced places in this comparison, including \(p=d\),
one can choose the class frames so that the substitution of \eqref{eq:unserved-normalized-frame} reads
\[
 Z_\tau\longmapsto \rho_\tau(P_{i+1})Z_\tau\rho_\tau(P_i)^{-1}.
\]
 The axis
generators and the resulting scalar factors in the coefficient
tests are units in the chosen coefficient valuation.
\end{proposition}

\begin{proof}
\proofstep{Choose simultaneous rational unitary frames.}

For each stage \(i\) and each of its classes choose a rational
unitary matrix \(P_i\) whose columns span its lattice at all
places above \(p\). Such a matrix exists simultaneously there.
The selected host extension \(E/F\) is unramified at finite places.
At a split place the two summands of a self-dual lattice are
mutually dual.
In paired local coordinates the Gram block between dual bases \(e_i,e'_j\) is
\[
  \big(h(e_i,e'_j)\big)_{ij}=I .
\]

At an unramified nonsplit place, orthogonal
splitting and surjectivity of the unramified unit norm identify
two unimodular Hermitian lattices. The resulting local
transporters are open. Multiplication by nearby scalar norm-one
elements avoids the eigenvalue \(-1\) at each place. Approximate
their inverse Cayley transforms in the rational vector space of
skew-Hermitian matrices by weak approximation over \(F\), and
apply the Cayley transform back. Thus both the columns and the
inverse coordinates clear into the required lattices by
integers prime to \(p\). The same matrix \(P\) works at every label over the place. Its inverse adjoint property is
\[
  P=(P^*)^{-1},\qquad P^*P=I.
\]

\proofstep{Separate the residual axes.}

For a CM type embedding \(\rho_\tau\) at \(w\) write \(d_\tau\)
for the image of the generator of its axis in the module
spanned by \((\rho_z a)_z\), \(a\) in the inverse half different.
At these enhanced places that ideal is the unit ideal locally.
The entries are integral. For two distinct type labels \(\tau,\tau'\), one can consequently separate their residue embeddings by an integer:
\[
  a\in{\cal O}_E,\qquad \overline{\rho_\tau(a)}\ne\overline{\rho_{\tau'}(a)}.
\]
 Indeed the labels
all extend the same embedding of \(K\), and the relative inertia
in \(E/K\) is trivial at these primes.

Distinct primes are
separated by the Chinese remainder theorem and distinct
embeddings of the same unramified local field are separated
on the residue field. This also applies at \(p=d\), where
\(K/\mathbb Q\) itself is ramified. Acting on the embedded
all-ones vector with the product, over \(z\ne\tau\), of
\(a_z-\rho_z(a_z)\), where
\(\rho_\tau(a_z)-\rho_z(a_z)\) is a unit, isolates the
\(\tau\)-axis with unit coordinate. The module is stable under
these actions. Its axis generator is therefore a unit, and
membership of the generator supplies a finite expression
\[
 d_\tau {\bf1}_\tau=\sum_l r_l(\rho_z a_l)_z,
 \qquad |r_l|_w\le1,\quad a_l\in\mathfrak b^{-1},
\]
where \(\mathfrak b\) denotes the half different.
This last expression keeps its global period directions.

\proofstep{Absorb the unit without changing the norm tests.}

Set \(a_\tau=d_\tau^{-1}\) and use the frame
\(a_\tau^i\rho_\tau(P_i)\) at stage \(i\). In the incident
Hom variable the factor
\(d_\tau a_\tau^{i+1}/a_\tau^i\) is exactly one. The substitution
on the series in \eqref{eq:unserved-reader} is now
\begin{equation}\label{eq:unserved-normalized-frame}
 Z_\tau\longmapsto
    \rho_\tau(P_{i+1}) Z_\tau\rho_\tau(P_i)^{-1},
\end{equation}
with the row of \(Z_\tau\) indexed by the target. These are
admissible choices for the true columns.

The product of the
column minors of a diagram \(\nu_\tau\) gains
\((a_\tau^i)^{|\nu_\tau|}\) on input frames. A determinant twist
of exponent \(s\) gains \((a_\tau^{iN_i})^s\), including signed
exponents. The target formulas use \(i+1\). Every such factor
is a unit, and homogeneity gives the same power on every term
of the corresponding polynomial reading. Thus the original
coefficient norm tests survive this change of frames, while the
Taylor map itself has the precise coefficient in
\eqref{eq:unserved-normalized-frame}. \label{cor:unserved-kernel-charts-2}

\end{proof}

\paragraph{One curve model before the inert orbit estimates.}
\begin{proposition}[One inert curve model for all translates]\label{prop:unserved-common-inert-model}
At the selected inert places \(p\ge5\) with \(p\ne d\), the two classes and all their prime-to-\(p\) translates use one good curve law, tangent unit \(c_w\) and quadratic correction \(q\). In good curve coordinates with logarithm \(l\), the parameter change for all these charts is
\[
 t(X)=l^{-1}(c_w L(X)),\qquad c_w\in{\cal O}_w^\times.
\]

In the even and signal cases the resulting normal form is
\textup{\eqref{eq:tag-i4}}. The odd case uses the doubled orbit and
the zero-step case uses the repeated-variable one-letter
phase. The phase and its coefficient estimates are uniform
in the labels, although the integral series may depend
on the classes and on every variable.
\end{proposition}

\begin{proof}
\proofstep{Fix the model and the common unit.}

The four cases treated here share the following hypotheses on the residue characteristic
\[
  p\ge5,\qquad p\ne d,\qquad \left(\frac{-d}{p}\right)=-1.
\]
Apply Proposition~\ref{prop:good-inert-cm-model} to the discriminant-normalized
period lattice. Changing the integral ideal representative of the
curve prime to \(p\) changes the discriminant root by the inverse
homothety and leaves its scaled lattice fixed. The good-model
homothety is a unit since its twelfth power compares two unit
discriminants. Write \(l\) for the good curve logarithm and \(L\) for the comparison logarithm. The integral parameter change \(t\) fixes the tangent unit through
\[
  l(t(X))=c_wL(X),\qquad t'(0)=c_w\in{\cal O}_w^\times.
\]
Use this one choice of parameter, Frobenius and curve phase \(q\) throughout the place. These precede the choice of the two classes, their ranks,
their diagrams, and the translations in the reader.

\proofstep{Choose the underlying Hom frame before the translates.}

Here is a rational orthonormal chart for the Hom of the actual
pair. For the standard lattices one matrix coordinate carries
the trace Hermitian form on the inverse half different over \(K\).
It is unimodular at \(p\) by duality of the trace lattice
(the half different times its conjugate is the relative
different).
Split off a vector of unit norm from its reduction.
For a lift \(u\) with unit norm, let \(\ell\) be its Hermitian pairing in the linear variable. The integral projection and its kernel are
\[
  x\longmapsto \ell(x)\ell(u)^{-1}u,\qquad \ell(u)\in{\cal O}^{\times},\qquad \ker\ell=u^\perp.
\]
 Such a
vector exists in every nonzero nonsingular Hermitian space over
the quadratic residue field, by polarization. Lift it to the
localized rational lattice.

Orthogonal projection onto its
complement divides only by that unit, so iteration gives an
orthogonal basis with positive rational norms \(n_j/e_j\) whose
numerators and denominators are prime to \(p\). Positivity and
rationality follow from definiteness and the fixed field of
conjugation on \(K\). In four copies of each line choose four
integer squares summing to \(n_je_j\). If \(Q_j\) is their quaternion matrix, then
\[
 Q_j^tQ_j=n_je_j I_4,\qquad
 (Q_j/n_j)^t\bigl((n_j/e_j)I_4\bigr)(Q_j/n_j)=I_4 .
\]
Both \(Q_j/n_j\) and its inverse are integral at \(p\). Use these four copies in each
\(2\)-by-\(2\) block of matrix entries, possible since both
ranks are even. Transport the resulting basis by

\[
  X\mapsto P_{i+1} X P_i^{-1}.
\]
 This is a rational isometry
and spans the Hom lattice over the local quadratic integers.
The basis is chosen before any translate.

\proofstep{Descend the corrected integral section from the fourth tensor.}

On four copies the normalized theta expansion of
Section~\ref{subsec:theta-normalization} has integral coefficients
in the locally matched frame. Its arguments are integral
unitary transforms of the orthonormal directions just
constructed. The one-variable translated section multiplied by
\(\exp(q)\) is integral in group coordinates. For the sum of
phases a unitary change preserves its class modulo logarithms
of principal integral units by the polarized CM norm relation
\eqref{eq:cm-norm-relation}.

Apply the change and extract four
equal coefficient indices on disjoint copies. This proves, in the
orthonormal basis \(b_j\), that the product
\[
  \exp\!\left(\sum_j q(X_j)\right)
  \Theta\!\left(\left(c_w\sum_j\rho_\tau(b_j)L(X_j)
                        \right)_\tau\right)
\]
is integral, for each of the normalized translates \(\Theta\)
in the finite reader. It holds also for an integral linear
combination since the multiplier and arguments are common. A pair
of lattices agreeing only locally at \(p\) gives the same conclusion.
Two integers prime to \(p\) clear the source and target lattice
inclusions in both directions, preserving these local arguments
and prime-to-\(p\) torsion.

\proofstep{Read the phases in the trace-dual cycle variables.}

In an even orbit expand the two parities in a basis and its
trace dual over the quadratic local field. For one pair of parity vectors \(e,o\) in a matrix coordinate, the Gram block is
\[
  \begin{pmatrix}h(e,e)&h(e,o)\\ h(o,e)&h(o,o)\end{pmatrix}=\begin{pmatrix}0&1\\1&0\end{pmatrix}.
\]
 This follows by multiplying the unitary
source and target changes in
\eqref{eq:unserved-normalized-frame} by their inverse adjoints
on the conjugate coordinates. Thus the change from the
orthonormal variables gives the hyperbolic substitution of
Section~\ref{subsec:curve-exponents}, with the trace-dual sums
there called \(E\) and \(O\). The root-of-unity average over
\(p^{d_p}-1\) elements is integral and replaces the summed
polarization by \(B_{d_p}\) modulo logarithms of integral
units.
The averaging denominator is a unit, since
\[
  p^{d_p}-1\equiv-1\pmod p,\qquad v_p\big((p^{d_p}-1)^{-1}\big)=0.
\]

Returning to additive coordinates gives \eqref{eq:tag-i4} of
Section~\ref{subsec:incident-sections}, with exponent
\(-B_{d_p}\) on each entry and argument scale given by the
common tangent unit. This works for any choice of the orbit
representatives and includes the signal orbits. The same
orthonormal identity at odd order is used below with the
trace-dual variables on the doubled orbit. At order one the
labelled basis is already unitary over the quadratic integers.
The norm relation transports the one-variable gauge, and its
polarization at the repeated variable gives the one-letter
phase with the factor \(1/2\). Thus in each case the integral
factor can depend on the classes and on every variable, while
the phase and its coefficient estimates are uniform in both
classes and all entries of their diagrams. \label{cor:unserved-kernel-charts-3}
\end{proof}


\subsubsection{Even cycles and signal primes}

On an even cycle the cyclic jet estimate gives a contribution from its paired positions. For a signal prime the same accounting must count each selected cycle only once.

\begin{proposition}[Entries on the even cycles]\label{prop:unserved-even-entries}
In the ordinary cutoff range of Lemma~\ref{hyp:unserved-level-free} with \(p\ge5\), suppose the
inert step has positive even order or the \(p\)-cycles are signal
cycles. In the latter case the cycles are counted once, with \(A\)
empty. For the full entry \(b\), using the even-cycle allowances below,
\[
 v_p(b)\ge\sum_r g_r+F_{\rm in}-E_{\rm I}.
\]
The estimate includes a signal cycle of length two.
\end{proposition}

\begin{proof}
\proofstep{The zero-index estimate on the first cycles.}

A positive even order of the inert Frobenius step gives one
set \(A\) of cycles on the first embedding and one set
\(BC\) on the second. At a signal prime Frobenius changes the
real embedding and the cycles are counted once, with
representatives on the first embedding. Then \(A\) is empty.
In both cases the cycle length \(d_p\) is even and at least
two, and \(p\ge5\). On a cycle in \(A\) the estimate with
index zero in \eqref{eq:tag-j1} implies the gain \(\max(g_0(r),0)\) at its
representative \(r\). In detail, the exclusion at a top
omits each least forward odd distance to a top, hence
the remaining terms in \(g_0(r)\) occur in the first sum of
\eqref{eq:tag-j1}. Its numerator \(\alpha_s\) is at most the diagram's
number of height-one columns and its recipient width
bound is at least the actual remaining width. With

\[
  \Lambda=\log_p(2a)\ge1,\qquad 1+\log_p(2a)\le 2\Lambda,\qquad 1+\log_p(a^2)\le3\Lambda,
\]
 the digit counts for a
height-one count at most \(2a\) and a remaining size at
most \(a^2\) are at most \(2\Lambda\) and \(3\Lambda\).
The error in \eqref{eq:tag-j1} is thus at most
\[
   5d_p M^2\Lambda+4\,\#\{\text{tops on }r\}\Lambda^2 .
\]
This is covered by the digit terms in \(g_0(r)\).
The nonnegative alternative in \eqref{eq:tag-j1} gives the maximum
with zero. This argument uses the zero-index coefficient
of the same homogeneous jet and is independent of the
choice of the neighbor summands on any other cycles.

\proofstep{Cyclic flow and the top capacity.}

For \(BC\) use the cyclic bound of
Proposition~\ref{prop:cyclic-jet} with the stage weights.
The capacity comparison \eqref{eq:flow-comparison} gives
the cyclic flow term 
\[
  F_r={\cal F}_p((\kappa w_z)_{z\in r})
\]
and the top contribution. At the second embedding there
are no tops and the top contribution is zero. On a signal
cycle the top contribution pays \(\max(g_0(r),0)\) after
the top error. Denote the resulting error by
\(\operatorname{err}_r\). Put \(n_r=\#\{\text{tops on }r\}\)
and \(W_*=b+\eta_{\rm sz}+1+3S_{\rm sz}\), the uniform
upper bound for the remaining width at a top. The error is
\[
 \begin{split}
 \operatorname{err}_r={}&\gamma n_r W_*+
  \frac{c_6}{p}\sum_{z\in r}w_z+
  d_p N_{j-1}N_j\bigl(1+\log_p(4d_p M a)\bigr)\\
 &+K_1d_p M^2\log_p(2a)+K_2 n_r\log_p^2(2a).
 \end{split}
\]
The constants satisfy \(c_6\ge12\), \(K_1\ge5\),
\(K_2\ge4\). The first two terms dominate
the top width and the loss \(11\sum w_z/p\) in
\eqref{eq:flow-comparison}.

The third and fourth terms cover
\(d_p N_{j-1}N_j(1+\log_p(a^2))\), since \(4d_pM\ge1\)
and \(K_1 M^2\log_p(2a)\ge N_{j-1}N_j\log_p a\).
We obtain the bound
\[
  \begin{cases}
    g_r, &r\in A,\\
    g_r+F_r-\operatorname{err}_r, &r\in BC,
  \end{cases}
\]
where \(g_r=\max(g_0(r),0)\) on first-embedding representatives
and is zero on second-embedding representatives. On the signal
cycles each odd step goes to the second embedding, which has
no tops. The distance in \(T_s\) is therefore infinite.
More explicitly, for any first-embedding top \(s\) on a signal cycle one has
\[
  \{j>0:j\ {\rm odd},\ s+j\text{ is a top}\}=\varnothing .
\]
The recipient width bound and the lower bound on the top's
height-one count show that \(T_s\) dominates the positive-part
sum in \(g_0\) term by term. That sum before the digit correction
is nonnegative. This pays \(g_r\) before the other errors in
\(\operatorname{err}_r\).

\proofstep{Add the bounds over the actual orbits.}

The logarithm and its inverse preserve disjoint sets of orbit
variables. The integral factor of the full theta series can
depend on all orbits. Expand it in monomials before inverting
the logarithm. Each monomial gives a product of admissible
single-orbit factors, and only finitely many meet a specified
jet. The derivative polynomial is a product over the orbits
on such a term.
For a monomial product \(\prod_r H_r(Z_r)\) in disjoint orbit variables the jet factors as
\[
  \partial^\alpha\!\left(\prod_r H_r\right)(0)=\prod_r\partial^{\alpha_r}H_r(0),\qquad\alpha=(\alpha_r)_r.
\]
 The valuation bounds therefore add. This
also explains the restrictions to a single orbit: setting the
other variables equal to zero commutes with their logarithms
and deletes their phase, whose constant is zero. The common
scaling of logarithmic coordinates is by a unit.

Each orbit term \(F_r\) is invariant under rotation, so
\(F_{\rm in}=\sum_r F_r\) is also the sum over all labels of
the same orbit terms divided by \(d_p\). The error
\(E_{\rm I}\) includes \(\sum_{BC}\operatorname{err}_r+
\max(0,\sum_A F_r)\). The summed entry estimate is therefore
at least \(\sum_r g_r+F_{\rm in}-E_{\rm I}\). In particular
the real-embedding count here is also valid for a signal
orbit of length two.

\end{proof}

\subsubsection{Odd cycle lengths}

An odd cycle has no alternating pairing of all its positions. Its contribution is estimated from the forced exponent and the ordinary error for this orbit.

\begin{proposition}[Entries on the odd cycles]\label{prop:unserved-odd-entries}
For \(p\ge5\) in the ordinary cutoff range of Lemma~\ref{hyp:unserved-level-free}, let \(b\) be the entry at a positive odd inert step of order \(d_p\ge3\). Then
\[
 v_p(b)\ge F_{\rm in}-E_{\rm I}.
\]
 There is no additional \(g_0\) gain
in this case. The phase comparison uses the doubled even cycle,
and the flow error includes \(\kappa\sum_{BC}w_z/p^3\).
\end{proposition}

\begin{proof}
\proofstep{Realize the phase on the doubled cycle.}

Here \(d_p\ge3\) is odd and the two sets of cycles again lie
over separate real embeddings. In the standard lattice
bases the local chart has an integral factor in the group
coordinates and an exponential of the sum of orbit phases
in additive coordinates. We explain the chart comparison
because the actual orbit has odd order and the parity in the
phase calculation is on a different set of letters.

Write \(\sigma\) for Frobenius on the unramified local
coefficients. Take an integral basis \(\omega_\rho\) of the
unramified extension of degree \(d_p\) over the unramified
quadratic field, orthonormal for the trace Hermitian form
whose conjugation is \(\sigma^{d_p}\).
The two parities belong to the doubled indexing set: reducing positions to the actual odd cycle identifies
\[
  s,\ s+d_p\ \in\mathbb Z/(2d_p)\quad\longmapsto\quad s\bmod d_p,\qquad (-1)^{s+d_p}=-(-1)^s.
\]
 This form is
unimodular by the unramified trace pairing. Reduction,
orthogonal splitting and surjectivity of the norm on
quadratic units give the basis, as in
Section~\ref{subsec:theta-normalization}. Its trace-dual basis
is \(\omega_\rho^*=\sigma^{d_p}\omega_\rho\).
In these coordinates trace orthogonality is
\[
  \operatorname{Tr}\big(\omega_\rho^*\omega_{\rho'}\big)=\delta_{\rho\rho'}.
\]

On a cycle with representative \(r\) and fixed matrix entry,
the trace-coordinate linear variables are
\[
 L_{r,\rho}(Z)=\sum_{\substack{0\le s<2d_p\\s\ {\rm even}}}
              \sigma^s(\omega_\rho^*) Z_{r+s}.
\]
Here \(r+s\) refers to the actual orbit of size \(d_p\).
The even labels give every position once. Integral standard
bases for the two incident lattices give integral Frobenius
coordinates \(g_s\) with 
\[
  \sigma^2 g_s=g_{s+2}.
\]
 Their
conjugate components give the paired evaluations. The
coefficients 
\[
  \sum_{s\ {\rm even}}\sigma^s(\omega_\rho)g_s
\]
are fixed by \(\sigma^2\), and the trace-dual identity gives
the inverse expansion of \(g_s\) in the displayed linear
variables. In the orthonormal kernel frame the source and
target evaluations are in perfect duality, entry by entry.
Expanding the Hermitian Gram matrix of this change of
variables and using trace duality thus gives identity:
it is zero between distinct orbits or entries, and is
\(\delta_{\rho\rho'}\) in each such block. It follows that
this linear change preserves the integral exponential
comparison in the orthogonal chart.

In group coordinates the backward steps of length \(2e\)
in the generalized logarithm give the same \(L_{r,\rho}\)
after applying the one-dimensional logarithm to the
integral formal sum with the indicated coefficients.
Write \(\mathcal G_{r,\rho}\) for that logarithm in the group coordinates. For the even representative \(s\), its coefficient is
\[
  [Z_{r+s}^{p^{2e}}]\,\mathcal G_{r,\rho}=p^{-e}\sigma^{s+2e}(\omega_\rho^*).
\]
 This is the Frobenius-twisted
logarithm of the formal sum by \eqref{eq:integral-log}.

The even and odd
linear forms in the doubled cycle agree on the two
identified copies when the odd form uses \(\omega_\rho\):
translation by \(d_p\) swaps their parity and applies
\(\sigma^{d_p}\) to the coefficient. The norm relation
therefore takes the sum of the diagonal phases to one
half of the averaged two-slot phase, modulo a phase
whose exponential is integral in group coordinates.
Here division by two comes from polarization, and
preserves this congruence because \(p\) is odd. In the
doubled cycle the averaged phase has the hypotheses of
Lemma~\ref{lem:two-differences}.

\proofstep{Pay the exceptional weights by the odd-step inequalities.}

For a nonnegative family \(b_z\) satisfying
\[
              p^q b_z+b_{z+q}\ge1\qquad(q>0\text{ odd})
\]
the phase at weights \(\rho+b_z\) has weight at least
\(\theta\). The exceptional monomials in the doubled
cycle project to two entries related by an odd step,
possibly to the same label. The displayed inequality pays
their missing unit of weight.

If both exponents are powers
of \(p\) with a common positive power, removing that
common power reduces the constraint to one of the displayed
inequalities. The odd step is measured before reducing labels
modulo \(d_p\). The inverse logarithm has weight at least
\(\rho\) in every coordinate, by the same induction as for
the even logarithm. It follows that the full rescaled
series has weight at least zero. Thus, including the
factorials, the valuation of a derivative of diagram size
\(H_z\) at every label is bounded below by
\[
          \sum_z H_z(\gamma-b_z)-
                    f N_{j-1}N_j(1+\log_p(a^2)).
\]
One can use this monomial estimate before applying the
integral coefficients of the framed minors.

\proofstep{Choose the dual weights on the two embeddings.}

On first-embedding cycles use \(b_z=1/(p+1)\). This is
admissible and gives a nonnegative contribution to the first
sum. On the second embedding choose an optimum dual for
the finite program with capacities \(\kappa w_z\) and
positive odd steps strictly less than \(d_p\).

Choose it in
\([0,1]\) and raise each coordinate to at least \(p^{-3}\).
This adds all omitted constraints, since any omitted odd
step is at least three. The cost on the original capacities
is at most \(\kappa\sum w_z/p^3\).
Writing \(b'_z=\max(b_z,p^{-3})\) for this modification, one obtains
\[
  \sum_z\kappa w_z(b'_z-b_z)\le p^{-3}\kappa\sum_z w_z .
\]
 A vertex at which the
original dual is at least \(2/p\) has saturated capacity
and no outgoing flow, by complementary slackness.
The sum of its weights over these vertices is at most
\(\sum w_z/p\), by the total source capacity of the flow.

The increase from \(\kappa w_z\) to \(H_z\) can therefore
reduce the lower bound by at most \(8\sum w_z/p\).
Indeed use \(b_z\le1\) on those vertices and
\(\gamma-b_z\ge-1/p\) elsewhere. The digit term is at most
\(3fM^2\log_p a\) since \(p\le a\). The resulting bound is
\[
 \sum_{r\in BC}F_r-
     \frac{\kappa}{p^3}\sum_{BC}w_z-
     \frac8p\sum_{BC}w_z-3fM^2\log_p a .
\]
Here \(E_{\rm I}\) contains the additional term
\(\kappa\sum_{BC}w_z/p^3\), as well as the other
allowances for these losses. There are no tops in \(BC\).
Subtracting \(\max(0,\sum_A F_r)\) covers the entire
first-embedding flow term. Thus this bound is at least
\(F_{\rm in}-E_{\rm I}\). The operator-free gain \(g_0\)
is zero in the assembled estimate at these primes.

\end{proof}

\subsubsection{The one-letter and ramified cases}

A zero inert step leaves one letter at each real label. This reduces its contribution to the one-letter derivative estimate and leaves the ramified primes to be treated separately.

\begin{proposition}[Entries at a zero inert step]\label{prop:unserved-zero-entries}
At a zero inert step with \(p\ge5\) and the ordinary prime range
and diagram sizes of Lemma~\ref{hyp:unserved-level-free}, there
is one letter on each real label. By the jet estimate of
Lemma~\ref{lem:one-letter-derivatives}, the entry \(b\) satisfies
\[
 v_p(b)\ge F_{\rm in}-E_{\rm C}.
\]

\end{proposition}

\begin{proof}
A zero inert step gives one letter at each real place.
Its local phase is the two-slot phase evaluated on
\((Z,\eta_z Z)\) and divided by two, with \(\eta_z\)
a unit. The one-letter calculation
\eqref{eq:one-letter-jet} applies separately to each label.
Expanding the common integral factor in group coordinates
as for the even orbits gives the estimate
\[
          \rho\sum_z H_z-f N_{j-1}N_j(1+\log_p(a^2)).
\]
In this case \(F_{\rm in}=\rho\kappa\sum_z w_z\).
The digit allowance in \(E_{\rm C}\) is
\(3 fM^2\log_p a\), and its remaining
\(\rho\sum_z w_z\) is nonnegative. This proves the
needed comparison with \(F_{\rm in}-E_{\rm C}\).

\end{proof}

\begin{proposition}[Entries at the discriminant prime]\label{prop:unserved-discriminant-entries}
At the discriminant prime \(p=d\), use the paired embeddings at
each real label, the level-free test and the normalized frame
\eqref{eq:unserved-normalized-frame}. Write \(b\) for the entry of the word \(u\). The counts below and the common budget give
\[
 v_p(b)\ge-\tfrac12\sum_t c_t(u)\ge F_{\rm D}-E_{\rm D}.
\]
 This comparison has no \(X''\) cutoff.
\end{proposition}

\begin{proof}
\proofstep{Apply the paired jet bound to the diagram.}

At the discriminant, let \(A_t^\pm\) count columns and
let \(a_t^-=A_t^--(A_t^-)'\) count columns of height one
on the first embedding. For the paired quadratic form of Section~\ref{subsec:cycle-jets}, the bound at the ramified prime is \[ -\frac12\sum_t c_t(u),\qquad c_t(u)=H_{(-,t)}-\mathbf1_{t\ {\rm top}}a_t^- +\mathbf1_{t\ {\rm top}}A_t^+ . \] The \(t\)-pairs are the two embeddings of the same real label, including at a top. On a generic
minus-frame change the plus directions change by its
inverse transpose. On the first row use the plus
column count, and on the other rows use the minus
degree, as in the ramified calculation following
Proposition~\ref{prop:cyclic-jet}. The quantity subtracted
from the first-row count is only the number of
height-one columns. The budget defines

\[
  E_{\rm D}=F_{\rm D}+\tfrac12\max_u\sum_t c_t(u).
\]
It follows that each word's entry has valuation at
least \(F_{\rm D}-E_{\rm D}\), with no use of \(X''\).

\proofstep{Transport the curve phase in the ramified formal coordinates.}

Here is the source of the paired shape used at the
discriminant. In the good curve model write \(Z=l(t)\) for
the formal logarithm and \(\phi(Z)=\mathrm{ph}(l^{-1}(Z))\)
for the one-variable phase in additive coordinates. The
formal group law and \(l'(t)\), with \(l'(0)=1\), are
integral. To see the curve congruences at this particular prime,
use the discriminant-one model as in
Proposition~\ref{prop:level-free-theta-jets}; \(d\equiv7\pmod8\) and

\[
  p=d\ge7,
\]
 so a good short equation is available over a finite
extension. The sigma addition calculation of
Section~\ref{subsec:curve-exponents} gives the parallelogram
identity modulo logarithms of principal integral units in the
curve parameter. Its proof cancels the entire \(s_2\) quadratic
and uses only good reduction. Polarization by the integral formal
half therefore gives biadditivity modulo these logarithms. For
\(a=1+\sqrt{-d}\), \(n=a\bar a=1+d\) is a unit. The isogeny
identity \eqref{eq:cm-isogeny} gives

\[
  q_E([a]t)\sim n q_E(t):
\]
 each nonzero kernel point has
prime-to-\(p\) order and integral affine coordinates, and the
leading factor \([a]t/(at)\) is a principal integral unit.
Polarize this identity, substitute \([a]^{-1}\) in the second
argument and use \(n[a]^{-1}=[\bar a]\). Subtract the adjoint
identity for the scalar one by biadditivity to obtain the
adjoint identity for \(\sqrt{-d}\). In particular this computation
does not require the inert unramified formal module at the
discriminant.

Let \(B(Z,W)=\phi(Z+W)-\phi(Z)-\phi(W)\), put
\(Z_i=l(t_i)\), and write
\[
  C=B(Z_1+Z_2,Z_3)-B(Z_1,Z_3)-B(Z_2,Z_3).
\]
The first-level congruence gives \(U=\exp(C)\) an integral
unit. Therefore
\[
 \partial_{t_1} C=U^{-1}\partial_{t_1}U
    \quad\text{in the integral formal series ring}.
\]
The other two derivatives of this series are integral.
In the expansion of \(C\), every summand other than
\(\phi(Z_1+Z_2+Z_3)\) omits at least one variable, so
\[
 \left.\partial_{t_1}\partial_{t_2}\partial_{t_3}C
             \right|_{t_2=t_3=0}
     =l'(t_1)\phi'''(l(t_1)).
\]
Here \(l'(0)=1\), and \(l'\) is an integral unit.
It follows that \(\phi'''(l(t))\) is integral.
Applying \(l'(t)^{-1}d/dt\) repeatedly and evaluating at zero
shows integrality of \(\phi^{(j)}(0)\) for every \(j\ge3\).
The congruence at multiplication by \(-1\) gives
\(2\phi'(0)\) integral.

Apply the same logarithmic-derivative argument in two
variables to the adjoint congruence for
\(B(\alpha Z,W)-B(Z,\bar\alpha W)\) with \(\alpha=\sqrt{-d}\)
and \(p=d\). Its mixed derivative at the origin is
\[
  (\alpha-\bar\alpha)\phi''(0)=2\sqrt{-d}\,S ,
              \qquad S=\phi''(0).
\]
This is integral, and 
\[
  v(2\sqrt{-d})=1/2
\]
 in the normalized
valuation. Hence 
\[
  v(S)\ge-1/2.
\]

\proofstep{Compute the quadratic phase in the orthogonal model.}

In the normalized orthogonal coordinates separate
\(\tfrac12 SZ^2\) from \(\phi(Z)\). The remaining series
has integral derivatives. So does the inverse logarithm.
Write \(l\) for the group logarithm and \(g=l^{-1}\). Its invariant differential has integral coefficient series with constant one, and differentiation gives
\[
  l'(g(Z))\,g'(Z)=1,\qquad g'(Z)=\big(l'(g(Z))\big)^{-1}.
\]
Differentiating the first identity at zero determines each next derivative of \(g\) from the lower derivatives, derivatives of \(l'\) and integer multiplicities, since \(l'(0)=1\).
The integral factor of the orthogonal chart therefore
retains integral derivatives in additive coordinates.
To compute its quadratic part on the selected matrices, use
the standardizing fourfold isometry of
Section~\ref{subsec:theta-normalization}. In the one-dimensional
coordinates the translated integral factor is supplied by
\eqref{eq:torsion-difference}. For \(Q\) modulo the period lattice the one-dimensional factor used in that calculation has the form
\[
  G_{jQ}(t)=e^{q_E(t)}(U_{jQ}\varsigma_*)(l(t)).
\]

The proof by reduction of the multiples \(jQ\) works also on the good curve at \(d\); the same gauge \(e^{q_E(t)}\) occurs for every \(j\). For a rational Hom direction
\(v\) with \(v\Gamma\) in the period lattice, the argument matrix
of the fourfold isometry has entries in the quadratic integers
localized at \(p\), by the multiplier-ring calculation of
Section~\ref{subsec:baseline-jets}. Thus the bilinear contraction
of two such argument columns is congruent modulo
\(\pi=\sqrt{-d}\) to their Hermitian contraction. Conjugating
one factor changes an integral quadratic-field entry by a
multiple of \(\pi\). Equivalently for two columns \(u,v\) in the orthogonal
coordinates,
\[
       u^{\mathsf T}v-u^*v=(u-\bar u)^{\mathsf T}v
                                 \equiv0\pmod{\pi}.
\]
The Hermitian contraction \(u^*v\) is the original trace
Hermitian value because the fourfold map is an isometry.
This computation also gives zero cross contractions modulo
\(\pi\) between disjoint copies. It remains to evaluate the
trace form on the directional tests in
\eqref{eq:unserved-normalized-frame}.

\proofstep{Separate the paired labels by integral selectors.}

Take a row of \(P_i^{-1}\) and a column of \(P_{i+1}\).
At label \(z\) multiply each by an integer of \(E\) that is
a unit there and clears the functional into the source dual
and the vector into the target. For a lattice \(M\), membership after localization gives exactly the clearing operation:
\[
  u\in M_{\mathfrak p}\quad\Longrightarrow\quad su\in M\text{ for some }s\in{\cal O}_E\setminus\mathfrak p.
\]
 This works on
any admissible pair agreeing with the selected pair locally.

If \(R,C\) denote the cleared row and column, \(a_l C R\) is a
global direction in the period module for every \(a_l\) from the
axis expression above. Evaluate that expression with coefficients
\(r_l/d_z\), then divide by the two clearing units in the
chosen embedding.
In fact the clearing factors \(s_1,s_2\) and the selected \(d_z\) are units, so
\[
  v_w\!\left(\frac{r_l}{s_1s_2d_z}\right)=v_w(r_l)\ge0 .
\]
 This is an integral combination giving exactly
one column of \eqref{eq:unserved-normalized-frame} and zero at
the other labels. All higher derivative remainders remain
integral after this linear substitution and products, by the
chain rule. Here and below \(d_z\) denotes the axis generator
at the type label of \(z\).

Expand the Hermitian trace on the two rank-one directions.
Each embedding contributes its row pairing times its column
pairing times \(a_l\overline{a}_{l'}\). The selector in the
first direction isolates the chosen embedding. For the second
direction the needed comparison is
\[
 \rho_{-,t}(\bar s)\equiv\rho_{+,t}(s)\pmod{\pi}
                  \qquad(s\in\mathcal O_E).
\]
Indeed the inertia involution at \(d\) has product labels
\(((1,-1),1)\). Composing the minus label with this involution
and complex conjugation gives the plus label at the same real
index \(t\). Inertia acts trivially on the residue field.
Its fixed part is unramified and the remaining quadratic field
is generated by \(\sqrt{-d}\). The extension over \(K\) is
unramified, so \(\pi\) remains a uniformizer.
Consequently the second selector kills all other partners
modulo \(\pi\). The source and target contractions on the surviving
pair give identity matrices modulo \(\pi\).

For instance if the
two row-clearing scalars are \(s_-,s_+\), the unspecialized
pairing is \(s_-\bar s_+\delta_{aa'}\), by the inverse-adjoint
identity for the two bases. After the two divisions the diagonal
entry is 
\[
  \rho_-(\bar s_+)/\rho_+(s_+)\equiv1\pmod\pi.
\]
The column contraction has the same calculation. Thus the
coefficient \(u_t\) in the surviving block is the scalar sum
of the products
\[
 \frac{r_{-,l}}{d_{-,t}}\,
 \rho_{-,t}(a_{-,l}\bar a_{+,l'})\,
 \frac{r_{+,l'}}{d_{+,t}},
\]
with the selectors for this pair of labels. Its terms are
integral because the inverse half different is integral here.
\proofstep{Extract the common quadratic phase from the fourth tensor.}

The opposite pairing gives the same coefficient modulo \(\pi\)
by symmetry of the original quadratic contraction. In particular
\(u_t\) can be fixed before the translates and their denominator
clearings. The quadratic sum on the four copies is therefore
the sum on each copy of
\[
        2\sum_{t,a,b}u_t X_{tab}Y_{tab}
\]
up to \(\pi\) times an integral quadratic polynomial. Because
\(v(S\pi)\ge0\) and two is a unit, the exponential of the error
multiplied by \(S/2\) has integral ordinary derivatives. This
follows also by the exponential product rule for derivatives
at zero.

Removing the common quadratic phase on the four
copies and taking the derivative at four equal indices bounds
the corresponding derivative on a single copy. Thus each
translate in \eqref{eq:unserved-reader} in its actual frame
is \(\exp(-S\sum u_tX_{tab}Y_{tab})\) times a series with
integral ordinary derivatives. The common exponential factors
out of the finite sum, including for the selected twist with
\(p\notin Q(\mu)\). This supplies the paired shape used in the
ramified bound. \label{cor:unserved-period-1}

\end{proof}

The base entry bound is zero at a level-free place away
from the discriminant whenever only integral jets are
used. This applies beyond the cutoff and at primes
with no applicable ordinary forced contribution. Section~\ref{subsec:baseline-jets} proves it in the actual
local bases, including at two and three. The stage bounds above are
uniform in the two classes of the entry, the word, and
each allowed tuple. Their inputs are the standard
theta charts of the same kernel datum, not charts
chosen separately for different coefficients of a jet.
\label{cor:unserved-period-2}


\subsection{From stage estimates to the column lattice}
\label{subsec:division-shift}

A bound for a matrix entry becomes a bound for a column only after the terminal restriction, the class lattice and its level trace have all been identified integrally.

\subsubsection{The integral restriction at the terminal classes}
\label{subsec:terminal-realization}

The last restriction in the chain is a branching map between interlacing types. Its matrix must be evaluated in the actual class frames before composing the finite kernels.

Fix a word and either the reference tuple or one of its used
twists. Write \(N=n+1\) for the last two ranks. At embedding
\(\tau\) the two true weights \(\lambda,\sigma\) interlace,
and \(\lambda_N=\sigma_n\). Indeed the row weight is the
interlacing row of the last Harish--Chandra parameter, and
the shift in both ranks is \(-m(t+1)/2\), with \(t\) odd.
We use the polynomial weights after subtracting the common
last part \(q\). The actions on full weights include the
characters \(\det^q\).

\begin{proposition}[Integral terminal branching]\label{prop:integral-terminal-branching}
For interlacing weights with the same last part, the
normalized polynomial lowering operator below intertwines
the branching spaces and carries the integral Weyl
lattice into the indicated integral row lattice on
the standard frames, simultaneously at all embeddings.
Writing \(q\) for the common last part, the same determinant twist restores the two full weights
\[
 q=\lambda_{n+1}=\sigma_n,\qquad \chi(g)=\det(g)^q.
\]

\end{proposition}

\begin{proof}
In a single matrix of variables, let \(\nabla_\eta\) be the
integral span of products of flag minors of diagram
\(\eta-\eta_{\rm last}\), and \(h_\eta\) the product with all
row sets initial. Use the Weyl lattice \(\Delta_\eta\) generated
from \(h_\eta\) by divided powers of row operators.

On polynomial
models in opposite frames their perfect pairing is
\[
 (f,g)_\eta=
 \frac{\sum_\nu \nu!\,[X^\nu]f\,[X^\nu]g}
      {\sum_\nu \nu!\,([X^\nu]h_\eta)^2}.
\]
This is integral Weyl--dual Weyl duality for the split general
linear group over \(\mathbb Z\), applied to the Schur module
of the indicated partition. The normalization follows by pairing
the two highest polynomial vectors to one. Equivalently the
semistandard bideterminants are an integral basis for the flag
span. Pairing them with the divided row words for the same
tableaux gives a unitriangular matrix in tableau order. This
is the standard monomial proof of duality and shows why extending
the coefficient ring preserves it. On several embeddings take
the tensor product of these pairings. \label{cor:terminal-realization-1}

Put \(a_i=\sigma_i-\lambda_{i+1}\), \(b_i=\lambda_i-\sigma_i\)
for \(1\le i\le n\). Let \(D_i\) be the initial \(i\)-minor,
and \(D_i^*\) the minor in the first \(i\) columns with rows
\(1,\ldots,i-1,N\). The vector
\[
              \xi=\prod_{i=1}^n D_i^{a_i}(D_i^*)^{b_i}
\]
is nonzero and is a highest vector of small weight \(\sigma\)
in the restriction of the large representation. This follows
from alternation in the first \(n\) rows and the row counts,
including the determinant character. The classical interlacing
branching rule in characteristic zero makes this summand simple
of multiplicity one. Normalize its embedding \(J_0\) by
\(J_0 h_\sigma=\xi\).

Here the integral map in the period has a specified scalar.
Define \(d_0=\sum_i a_i\) and
\[
 R_i=[\lambda_{i+1},\min(\lambda_i,\lambda_{i+1}+d_0)]\cap\mathbb Z,
 \qquad
 z(\nu)=\prod_{i=1}^n\prod_{v\in R_i\setminus\{\sigma_i\}}
                      \prod_{j=1}^n(v-i-\nu_j+j).
\]
Let \(c=z(\sigma)\). The shifted interlacing intervals
\([\lambda_{i+1}-i,\lambda_i-i]\) are pairwise disjoint.
Thus \(c\ne0\). If \(\nu\) interlaces and
\(\sum\nu_i=\sum\sigma_i\), then every \(\nu_i\) belongs to
\(R_i\). Unless \(\nu=\sigma\), substituting \(v=\nu_i\) for
a differing coordinate gives \(z(\nu)=0\). The Capelli
characteristic polynomial for the small row action gives a
central operator with eigenvalue 
\[
  \prod_j(x-\nu_j+j)
\]
on highest weight \(\nu\).

In the usual column determinant
convention evaluate the Capelli determinant of
\(E+\operatorname{diag}(n-1,\ldots,0)-uI\) at \(u=x+n\)
and multiply by \((-1)^n\). The eigenvalue is the one just
stated. \cite{molev-branching} The row Lie operators \(E_{ij}\), with the integral
determinant twist on the diagonal, preserve the integral flag
span by differentiation of the integral action. Taking the
product at \(x=v-i\) gives an integral operator \(Z\) acting
by \(z(\nu)\). On the subspace of last-coordinate weight
\(\sum\lambda_i-\sum\sigma_i\), this means
\[
       Z|_{V_\nu}=z(\nu)\,\mathrm{id}_{V_\nu}
               =c\,\mathbf1_{\nu=\sigma}\,\mathrm{id}_{V_\nu}.
\]
The multiplicity-one summand \(V_\sigma\) is generated by
\(\xi\). Thus \(Z\) is \(c\) times that projection.

We verify that \(J=cJ_0\) takes \(\nabla_\sigma\) to
\(\nabla_\lambda\). Label the factors in a product \(P\) of
small flag minors. For each \(i\ge2\) append row \(N\) to
\(b_i\) columns of height \(i-1\) in the tableau for \(P\).
There are at least that many since

\[
  \sigma_{i-1}-\sigma_i\ge\lambda_i-\sigma_i.
\]
Use \(b_1\) additional singleton columns with entry \(N\).
For each enlarged minor use the first \(i\) matrix columns,
and keep the other minors. This gives an integral polynomial
\(M(P)\) of large shape and the required last weight. The
choice of which labeled columns to enlarge is fixed as the
row indices vary. The initial product maps to \(\xi\).
Cauchy--Binet shows that applying any small elementary row
substitution to \(P\) and \(M(P)\) gives expansions with the
same coefficients indexed by the same changes of the small
row tuples. The appended row is fixed. This remains true for
each coefficient in the substitution parameter and hence for
all sequences in divided row operators.

For each tuple let \(D(P)\) be the error term in the \(\sigma\)-summand. The substitution preserves the small diagonal torus weight, so whenever the difference is nonzero it satisfies
\[
  D(P)=ZM(P)-cJ_0P,\qquad \operatorname{wt}(D(P))=\operatorname{wt}(P).
\]
If this weight differs from the highest weight, its Fock pairing
with \(\xi\) vanishes by torus invariance of the Hermitian
version of the pairing.

At highest weight the row counts force
every nonzero small column to use rows \(1,\ldots,k\) for its
height \(k\). Count row 1 first, then row 2 in all columns
of height at least two, and continue. The signs of their
permutations agree after appending \(N\); a repeated index
kills both products. Hence at this weight the difference
vanishes. Put \(D(P)=ZM(P)-cJ_0P\). The simultaneous expansions give,
for every word \(T\) in the divided row operators,
\[
              0=\langle T D(P),\xi\rangle
               =\langle D(P),T^*\xi\rangle .
\]
Adjunction interchanges the two row indices. The vectors
\(T^*\xi\) span the summand containing \(D(P)\), so \(D(P)=0\). This proves
integrality without requiring the unprojected lift \(M\) to
respect the straightening relations.

For instance at \(\lambda=(2,1,0)\), \(\sigma=(2,0)\) one
has \(c=4\). Write \(x_i=X_{i1}\) and \(D_{ij}\) for the
two-column minor on rows \(i,j\). Then
\(J_0(x_1^2)=x_1D_{13}\), and applying the row derivation
from 1 to 2 gives
\[
 J_0(x_1x_2)=(x_2D_{13}+x_1D_{23})/2 .
\]
Its numerator has odd coefficients. The unscaled embedding
already fails to preserve the flag span at two in this example.

Take the normalized transpose of \(J\). It maps
\(\Delta_\lambda\) into \(\Delta_\sigma\) by the perfect
pairing. For the coefficient extraction, set the last column to zero and take the operator
\[
  P\longmapsto\Big[\prod_i X_{Ni}^{b_i}\Big]\big(P\big|_{\text{last column}=0}\big).
\]
 On \(\xi\) the result is
\(h_\sigma\). Starting with the largest \(i\), the required
power of \(X_{Ni}\) uses all the diagonal last-row terms of
the remaining minors of height \(i\). Project the resulting
polynomial orthogonally to the small representation in the
Fock form, and multiply by \(c(\xi,\xi)_\lambda\).
This is the formula for \(\mathrm{br}_{\rm std}\). Coefficient
extraction commutes with the small row action; the projection
commutes by unitary invariance and complexification. Both
intertwiners agree on \(\xi\), where the value is
\(c(\xi,\xi)_\lambda h_\sigma\), so multiplicity one makes
them equal. All coefficients descend to the rational
polynomial model. Tensoring proves integrality in the
standard frames simultaneously at all embeddings.
\label{cor:terminal-realization-2}
\end{proof}

\paragraph{Lowering in the class frame.}

Restriction on highest-weight spaces is rational. Integrality of its values must be checked in the frames of the two matched terminal classes.

\begin{lemma}[Integral terminal values in the class frame]
\label{lem:terminal-integral-class-frame}
At a lifted class and its matched terminal class, write \(L_{\rm st}\) for the integral coefficient lattice in the stage frame and \(L_{\rm sh}\) for the sharp point lattice. For the last-line generator \(t_z\) and the central weight differences \(s_\tau=|\lambda_\tau|-|\sigma_\tau|\), the normalized restriction satisfies
\[
 {\rm br}_{t_z}=\left(\prod_\tau\tau(t_z)^{-s_\tau}\right){\rm br}_{\rm std},\qquad {\rm br}_{t_z}(L_{\rm st})\subseteq L_{\rm sh}.
\]
 It is covariant in the terminal class and commutes
with the reference and twisted evaluations at the respective tuples.
\end{lemma}

\begin{proof}
Use Proposition~\ref{prop:integral-terminal-branching} in the
standard coordinates before transporting the class bases.
At a terminal class let \(L'\) be its sharp lattice. In the
lifted class the decorated lattice is
\[
  \widetilde L=\{(v',v_N): v'\in L',\quad t_zv_N\in O_E\}.
\]
Here \(t_z\) is the last-coordinate multiplier of the
representative. Fix a rational prime \(p\), and choose a basis
matrix \(B'\) of \(L'\) at the places over \(p\). A single
matrix over \(E\) suffices, since the localization of the
Dedekind domain at this finite set is semilocal and its
fractional ideals are principal. A basis matrix of
\(\widetilde L\) there is \(B=\operatorname{diag}(B',t_z^{-1})\).

Let \(s_\tau=\sum_i\lambda_{\tau i}-\sum_i\sigma_{\tau i}\)
using the full weights. The actual restriction at the
representative uses
\[
   \mathrm{br}_{t_z}=
        \Big(\prod_\tau\tau(t_z)^{-s_\tau}\Big)\mathrm{br}_{\rm std}.
\]
Indeed the last diagonal group has character \(a^{s_\tau}\)
on the selected summand by the two central weights. In the
inverse-transpose representation \(\rho^\vee\), reading in
the inverse basis uses \(B^t\) on the polynomial variables.
The block identity is therefore
\begin{equation}\label{eq:terminal-block-identity}
 \rho^\vee_\sigma((B')^{-1})\,\mathrm{br}_{t_z}(G)
       =\mathrm{br}_{\rm std}\big(\rho^\vee_\lambda(B^{-1})G\big).
\end{equation}
In particular the last diagonal entry acts with the full
weight difference, even when the polynomial degree has
been reduced by a determinant twist.

At the lifted class evaluate an integral stage form. A change
from its chosen representative acts on its value by the same
lattice transport as on its decorated lattice, times the finite
character, a unit. Thus its coordinates in \(B\) belong to
the local Weyl lattice. Write \(F(z)\) for the value at the lifted class and
\(T(z')\) for its restriction at the matched terminal class.
By the stated restriction on point values,
\[
   \rho^\vee_\sigma((B')^{-1})T(z')
      =\mathrm{br}_{\rm std}\bigl(\rho^\vee_\lambda(B^{-1})F(z)\bigr).
\]
The right side lies in the standard integral Weyl lattice
by integral lowering. The finite character of the restriction
is the quotient of the two class characters. The same block
intertwining identity with the stabilizer action gives covariance
under a change of representative. Equivalently this is the
restriction of the complex automorphic form after embedding
the coefficients, and injectivity of the embedding gives the
covariance over \(K\).

A covariant family on a finite set of class representatives
with values in their \(K\)-spaces defines the rational terminal
point. Its localization is in the integral point lattice
precisely when all these values are locally integral.
For finitely many local values \(y_c\) with clearing factors \(s_c\) away from \(w\), a single factor works, since
\[
  s=\prod_c s_c,\qquad v_w(s)=0,\qquad sy_c=(s/s_c)(s_c y_c).
\]
 To
see the reverse inclusion, clear their finitely many
denominators by a single element outside the place. This
preserves the linear covariance equations. Intersecting over
the places gives the assertion for integral global inputs.
These comparisons use the lifted classes and \(t_z\) of the
tuple in question. They commute with reference or twisted
evaluation on that tuple. \label{cor:terminal-realization-3}
\end{proof}

\paragraph{Realizing a sharp test.}
The selected row uses the coordinate image of the integral
level lattice. Here is the descent to evaluations in that
lattice. Let \(j\) be the semilinear embedding of level
coordinates in the sharp complex row ambient space at this
tuple. Restricting an algebraic level form to class
representatives identifies it with its coordinates satisfying
the stabilizer equations. Conversely, if \(a\) is a chosen representative with
assigned value \(y_a\), the covariance operator \(R(g)\)
assigns \(R(g)y_a\) to its translate.

When \(ga=g'a\),
the stabilizer equation \(R((g')^{-1}g)y_a=y_a\)
makes these assignments agree.
Indeed covariance gives
\[
  R(g)y_a=R(g')R((g')^{-1}g)y_a=R(g')y_a .
\]
 This produces exactly the
algebraic form with those coordinates. The integral coordinate lattice
consists of the vectors whose values lie in the transported
flag spans at the level points. This description also results
if one first uses the full finite level action and then
restricts to the given character: a vector already in the
character space is recovered from its level coordinates by
inclusion.
In particular, for a character \(\chi\) of the level group, the integral points of its character space are given in level coordinates by
\[
  M[\chi]=\{v\in M:R(h)v=\chi(h)v\ \text{for every level element }h\}.
\]
 No projector denominator enters this operation.

Choose the designated finite generators \(v\) of the sharp
ambient lattice. In the selected row data they are the
embedded coordinate generators of the integral level lattice.
More explicitly, each has a coordinate lift \(g_v\) with
\(j(g_v)=v\), and a level coordinate \(y\) is integral exactly
when \(j(y)\) is an \(\mathcal O_K\)-combination of these
generators under the coefficient embedding. This description
follows by taking a finite generating set of the integral
coordinate module before extension of coefficients, and
transporting its span with the chosen level isomorphism.
If \(F=\sum_v\iota(c_v)v\) is a sharp integral test, take

\[
  y=\sum_v c_vg_v,
\]
 with the coefficients mapped to
\(\mathcal O_K\) if the record uses its original coefficient ring.
By semilinearity,
\[
          j(y)=\sum_v \iota(c_v)j(g_v)
                     =\sum_v\iota(c_v)v=F.
\]
Each generator \(g_v\) is in the integral level module and
so is \(y\). The stabilizer equations are linear, so the
covariant extension from \(y\) is integral at every class.
This uses the image of the integral module, including its
index in the rational coordinates.

Read that form at the terminal representative using
Lemma~\ref{lem:terminal-integral-class-frame}. By the decorated
lattice identification its value is in the flag span for
\(L'\). The sharp ambient identification and the inverse
period identification give the same function there: both
extend the same level coordinates by covariance.
Indeed if two covariant extensions \(F_1,F_2\) agree at the representative \(a\), then on its entire class
\[
  F_1(ga)=R(g)F_1(a)=R(g)F_2(a)=F_2(ga).
\]
 This applies
to the reference identification and to the twisted one with
its own representative. Thus \emph{every} sharp test is
represented at every terminal class by an integral row
polynomial for the lattice dual to the column just
constructed. The pairing in the terminal sum can now be
applied class by class. \label{cor:terminal-realization-4}

\hypertarget{source-terminal-period}{}
Use the column lattice \(N_u\) and the integrality tests on the row lattice \(M_\mu\) of Section~\ref{subsec:selected-period}. We describe
the effect of the class sums on the column and the terminal
reading on these tests. These steps take place over the
coefficient field \(K\). Write \(\mathcal O_w\) for its local
ring and 
\[
  v=v_w/e_w
\]
 normalized by \(v(p)=1\). Let \(D_{\Gamma,u}\) be the
common stabilizer multiple of the chain and put
\(\Delta_w=v_p(D_{\Gamma,u})\). The multiple includes the
terminal classes of the reference and used twisted chains.

\begin{proposition}[The terminal valuation without a neighbor operator]\label{prop:unserved-terminal-valuation}
At a level-free place without the local orbit condition,
with \(p\notin Q(\mu(w))\), suppose the stage coefficients
have normalized valuations at least \(B_j\) in the actual
incoming frames. For each terminal period value \(b\) on the local tests of \(M_\mu\),
\[
 v_p(b)\ge\sum_{j=1}^m B_j-m\Delta_w,\qquad \Delta_w=v_p(D_{\Gamma,u}),
\]
where \(D_{\Gamma,u}\) is the common
chain stabilizer multiple including both terminal class sets.
\end{proposition}

\begin{proof}
\proofstep{Pass from coefficient valuations to outgoing integral Weyl tests.}
At a level-free place the localized lattice of stage \(i\)
is read by the local Weyl values at each of its classes.
In the coordinates of the incoming member of the chain these
values are integral for every flag-minor test. Conversely,
for a vector in the rational stage space, integrality of
these values is the defining local lattice condition.

For
one incoming and one outgoing class the kernel reading is
the derivative polynomial in the incoming frame, multiplied
by a unit. Expand it in products of flag minors on the
incoming frame. Integral straightening as in
Section~\ref{subsec:cycle-jets} shows that a coefficient
bound \(B_j\) gives coefficients of valuation at least \(B_j\)
also in such a reading.
Each straightened coefficient is an integral linear combination; explicitly,
\[
  v(a_i)\ge B_j,\quad u_i\in{\cal O}_w\quad\Longrightarrow\quad v\!\left(\sum_i u_i a_i\right)\ge B_j .
\]
 Applying the incoming local
functional gives the same bound on the outgoing values.
The class map is rational and the stage space is stable
under the full operator. These two facts let one use the
local lattice criterion after rescaling the output. Indeed
a nonzero rational value of valuation at least \(B_j\)
has 
\[
  v_w(y)=e_w v(y)\ \ge\ \lceil e_w B_j\rceil\qquad(y\ne0,\ v(y)\ge B_j),
\]
 so the rescaling can be
by a power of a uniformizer chosen in \(K\). This applies
also when \(B_j<0\).

\proofstep{Charge an incoming stabilizer at each stage where it occurs.}
The full map is the sum of the class maps with weights
\(1/|\Gamma_c|\) on the incoming class \(c\). This costs
at most \(\Delta_w\) in \(v\) at stage \(j\) when \(j>1\).
The first incoming group is trivial for the reference
and every used twist. Consequently the accumulated
valuation of the stage-\(m\) vector, relative to its
localized stage lattice, is at least
\[
           \sum_{j=1}^m B_j-(m-1)\Delta_w .
\]
The same composition applies to a bound proved on
the entire column at a specified level with its specified
starting lattice. Products of the successive scalars add
their valuations, so composition introduces no rounding
loss. The estimates without the local orbit condition start in the ramified
submodule of the original trace-zero column. The
calculation at the prime assigned to the sequence uses the lattice with its extra
local condition at the assigned prime.

For a vector \(y\) in the localized stage-\(m\) lattice,
first take its image under the point restriction and the
lowering map into the last compact row model. The point
lattice comparison sends integral stage vectors into the
designated integral terminal lattice. For the reference
chain, the normalized Petersson period decomposes over
terminal classes as
\begin{equation}\label{eq:terminal-period-valuation}
     P_\mu(y,F)=
         \sum_{c}\frac{q_{\mu,c}(y,F)}{|\Gamma_{\mu,c}|}.
\end{equation}
The same formula holds for the used twist with its own
terminal class restrictions. It follows by restricting
both functions to each orbit in the finite double quotient.
The characteristic projectors of the classes are disjoint.

In the invariant Hermitian pairing each class contributes
the local pairing with weight \(1/|\Gamma_{\mu,c}|\).
Writing \(P_c\) for the period restricted to class \(c\) and \(\langle\ ,\ \rangle_c\) for its local pairing, this weighting reads
\[
  P_c(y,F)=|\Gamma_{\mu,c}|^{-1}\langle y_c,F_c\rangle_c .
\]
The real type pairing uses the orthonormal frame for the
positive terminal form and the normalized invariant product
in that frame. The same normalization is used for the
Weyl tests in the terminal lattice.

The integral row condition represents the restriction of
each sharp test \(F\) by a flag-minor polynomial on the
local terminal lattice at \(c\). This is the same lattice
whose derivative dual reads the point image of the
column. More explicitly, in a lattice frame take the
span of the flag-minor products of the prescribed diagram.
A change of basis of the underlying lattice and the
inverse-transpose change in its dual transport the
constant-coefficient derivative pairing. They carry
the row span and the dual column lattice together.
In the reference frame the numerator of the pairing for a row monomial \(X^\alpha\)
is given by \(\partial^\alpha\) at zero.

For monomials this numerator has the diagonal form
\[
  (\partial^\alpha X^\beta)(0)=\begin{cases}\alpha!&\alpha=\beta,\\0&\alpha\ne\beta.\end{cases}
\]
 Divide by the product
of the squared Fock norms of the highest polynomials, as in
Section~\ref{subsec:terminal-realization}. The distinguished
real normalized form is this polynomial pairing on the
restricted representation, with the row conjugation in
its coordinates. Thus evaluating an integral point
column on the integral row span gives the embedded
value of its \(K\)-valued Weyl pairing. Equivalently
multiply the period of the two class restrictions by
\(|\Gamma_{\mu,c}|\) to recover that pairing, since the
single-class weight in the real period is its reciprocal.

In particular, if \(P_c\) denotes the period of the two restrictions to this class,
\[
  q_{\mu,c}(y,F)=|\Gamma_{\mu,c}|P_c .
\]
This gives \(q_{\mu,c}(y,F)\) in \(K\) of valuation
at least zero. The same duality applies in the used
twisted chain with its transported row span. Its
integral period reading is used at \(p\notin Q(\mu)\).
The claim holds for localized \(y\) by \(K\)-linearity in
the column. For the normalized integral column, the scalar
\(D_{\Gamma,u}/|\Gamma_{\mu,c}|\) is integral for every class \(c\).
Consequently \eqref{eq:terminal-period-valuation} gives
\[
 D_{\Gamma,u}P_\mu(y,F)
   =\sum_c \frac{D_{\Gamma,u}}{|\Gamma_{\mu,c}|}\,q_{\mu,c}(y,F)
        \in\mathcal O_w,\qquad v(P_\mu(y,F))\ge-\Delta_w .
\]
This incurs one stabilizer loss for the entire terminal sum. A scalar removed from the column before this
reading multiplies the period and adds its own valuation.

It follows for the case without the local orbit condition that the valuation of the
rational period is at least
\begin{equation}\label{eq:unserved-terminal-bound}
                   \sum_{j=1}^m B_j-m\Delta_w .
\end{equation}
For a place with the local orbit condition the same terminal loss removes the
single \(+\Delta_w\) in the column bound. At the prime assigned to the sequence the stabilizer multiple is prime to \(T_u\), so the
terminal reading preserves the bound at the assigned prime.
Indeed the terminal denominator at this prime has
\[
  T_u\nmid D_{\Gamma,u}\quad\Longrightarrow\quad v_{T_u}(D_{\Gamma,u})=0 .
\]
 At
either type prime the local transport estimate is
already a bound on the period against the row test.

To apply this to the actual test space for \(w\), take
\(\mu=\mu(w)\). The level exclusion \(p\notin Q(\mu(w))\)
refers to the coefficient place of this reading. The row space at the refined level image of these tests, localized at \(w\), is \(R_w\).
The selected period comparison over \(K\) induces the same
valuation on the common rational pairing. The local bound
therefore applies to every test vector in \(R_w\), including
localized linear combinations of the images.

Section~\ref{subsec:terminal-realization} identifies these
classwise row spaces at the refined levels and terminal points on both tuples.
\label{cor:terminal-period-1}
\end{proof}



Let \(s_w\) be the bound obtained by summing the
placewise stage prices for the word \(u\). At an
ordinary place the prices are the gain plus \(F-E\),
with the class denominators included by
\eqref{eq:unserved-terminal-bound} of Proposition~\ref{prop:unserved-terminal-valuation}
or its version with the local orbit condition.
The selected lattice gain in \(s_w\) is bounded by the
gain proved for the local chain. At the discriminant the valuation bound is \(F_{\rm D}-E_{\rm D}-\Delta_w\), and at the
prime assigned to the sequence it is \(-E_T(w)\). The type-prime estimate
supplies \(s_w\) directly. Set \(\sigma_w=e_ws_w\).

Write \(L=K\) for the coefficient field. Let \(A_t\) be the row
quotient map followed by its rational identification with the
common row space. The common row lattice \(R\) is specified
by
\[
       R_w=(A_{t(w)}M_{t(w)})_w
\]
at the localized valuation rings. Suppose that the corresponding
bilinear form \(B\) satisfies \(P_t(\varphi,F)
=\xi_t B(\varphi,A_tF)\), with \(v_w(\xi_{t(w)})=0\), as in
the period comparison.

\begin{construction}
\begin{tabular}{@{}lll@{}}\toprule
 & Domain of the pairing & Ideal containing the values \\ \midrule
At \(w\) & \(N_w\times R_w\) & \(\mathfrak p_w^{\tau_w}\mathcal O_{L,w}\) \\
After twisting & \(I_wN_w\times R_w\) & \(\mathcal O_{L,w}\) \\
Globally & \(IN\times R\) & \(\mathcal O_L\) \\ \bottomrule
\end{tabular}
\par\smallskip
\(\tau_w=\lfloor\sigma_w\rfloor,\quad
 I_w=\mathfrak p_w^{-\tau_w}\mathcal O_{L,w},\quad
 I=\prod_w\mathfrak p_w^{-\tau_w}.\)
\caption{From the local period valuations to the integral Hermitian dual in Proposition~\ref{prop:finite-period-shift}.}\label{fig:analytic-shift}
\end{construction}

\begin{proposition}[The fractional-ideal shift]\label{prop:finite-period-shift}
Let \(N\) be the original column lattice of positive rank
with its Hermitian metrics, and let \(R_w=(A_{t(w)}M_{t(w)})_w\)
be the localized row lattice transported to the common rational space. Suppose the placewise period estimates give
\(v_w(B(N_w,R_w))\ge\sigma_w\) at each finite place \(w\), where
\(\sigma_w=0\) outside a finite set. Put
\[
 \tau_w=\lfloor\sigma_w\rfloor,\qquad
 I=\prod_w\mathfrak p_w^{-\tau_w},\qquad N'=IN .
\]
Then \(B(N')\subseteq R^\vee\), and the squared-norm slope of
\(N'\) differs from that of \(N\) by
\[
 2\sum_w \frac{\log N\mathfrak p_w}{[L:\Q]}\,\tau_w.
\]
Replacing \(\tau_w\) by \(\sigma_w\) in this sum has error at most
\(2\sum_{p\in S}\log p\) when all nonintegral \(\sigma_w\) are
above the rational primes in \(S\).
\end{proposition}

\begin{proof}
\proofstep{Check the inclusion at each finite place.}
Localizing identifies \(I_w\) with a principal fractional ideal
generated by an element of valuation \(-\tau_w\). Multiplying the
column by this ideal preserves integrality of \(B\) on \(R_w\),
because \(v_w(B)\ge\sigma_w\) for the unshifted generators.
In a uniformizer \(\varpi_w\), the local test on a nonzero generator reduces to
\[
  v_w(B)\in\mathbb Z,\quad v_w(B)\ge\sigma_w\quad\Longrightarrow\quad B\in\varpi_w^{\lceil\sigma_w\rceil}{\cal O}_w .
\]
Localization permits division by any scalar of valuation zero.
Intersecting over \(w\) now gives \(B(N',R)\subseteq{\cal O}_L\),
or \(B(N')\subseteq R^\vee\). The conclusion in fact holds with
any integral \(\tau_w\le\lceil\sigma_w\rceil\), by integrality
of the valuation at a nonzero coefficient.

\proofstep{Compare determinants and bound rounding.}
For the squared Hermitian slope used here put
\(\ell_w=\log N\mathfrak p_w/[L:\mathbb Q]\).
At rank \(n>0\) the \(w\)-adic determinant exponent of the
lattice changes by \(-n\tau_w\). The archimedean metrics are
fixed. The normalized log covolume changes by
\(-2n\sum_w\ell_w\tau_w\), and hence
\[
           \mu(N')=\mu(N)+2\sum_w\ell_w\tau_w .
\]
If the nonintegral \(\sigma_w\)'s occur only over a finite set
\(S\) of rational primes, replacing \(\tau_w\) in this last
formula by \(\sigma_w\) costs at most \(2\sum_{p\in S}\log p\).
Indeed \(0\le\sigma_w-\tau_w<1\) and
\(\sum_{w\mid p}e_w\ell_w=\log p\) by the degree formula. Both
the inclusion in the dual and the slope change use the very
same tuple selection at the finite places.
\end{proof}

\paragraph{The finite support in the selected family.}
The nonzero exponents in the sum have residue
characteristic among the primes at most \(X''\), the
primes \(T_u,g_0,g_1,d\), and the prime divisors of
\(D_{\Gamma,u}\). Above \(X''\) and outside those exceptions, the family prescribes zero gain. The forced term vanishes by its cutoff. The local error on the complement of the word, type and discriminant places has a residue-prime contribution only for \(p\le X''\). This leaves just \(v_p(D_{\Gamma,u})=0\)
at every lift. These identities give zero for the prescribed
shift. Integrality of the coefficient there makes this price valid. The product \(I\) is
therefore finite. The bound \(\vartheta(X'')\le X''\log4\)
and \(\sum_{p\mid D_{\Gamma,u}}\log p\le\log D_{\Gamma,u}\)
give the rounding allowance
\[
  2\bigl(X''\log4+\log(T_u g_0g_1d D_{\Gamma,u})\bigr).
\]

This rounding allowance is independent of the number of coefficient places
and the degree of \(L\).
For the bound on \(\vartheta\), use the divisibility of
\(\binom{2n}{n}\) by every prime in \((n,2n]\).
For \(2n+1\), the analogous primes in \((n+1,2n+1]\)
divide \(\binom{2n+1}{n}\le2^{2n}\), by the two equal
central binomial coefficients. Induction on the integer
argument gives \(\vartheta(x)\le x\log4\).

Thus Proposition~\ref{prop:finite-period-shift} applies to the
selected tuples with a rounding allowance independent of the
degree of the coefficient field. The shift can therefore be used in Section~\ref{sec:word-division} on the original column, even when the local coefficient field had to be enlarged to compute a kernel entry.
\label{cor:finite-valuations-1}
